% Task56 interim draft. Generated tables and figures are the numerical source.
% Default: anonymous review. Use \def\AuthorDraft{1}% Task56 interim draft. Generated tables and figures are the numerical source.
% Default: anonymous review. Use \def\AuthorDraft{1}% Task56 interim draft. Generated tables and figures are the numerical source.
% Default: anonymous review. Use \def\AuthorDraft{1}% Task56 interim draft. Generated tables and figures are the numerical source.
% Default: anonymous review. Use \def\AuthorDraft{1}\input{main.tex} for placeholders.
\newif\ifshowauthors
\ifdefined\AuthorDraft\showauthorstrue\else\showauthorsfalse\fi
\ifshowauthors
  \documentclass[sigconf,review]{acmart}
\else
  \documentclass[sigconf,review,anonymous]{acmart}
\fi
\setkeys{acmart.cls}{balance=false}
\setcopyright{none}
\settopmatter{printacmref=false}
\acmConference[Draft]{Unsubmitted research draft}{November 2026}{No venue assigned}
\acmDOI{}
\acmISBN{}
\usepackage{amsmath}
\usepackage{graphicx}
\usepackage{booktabs}
\newcommand{\sys}{\textsc{ckks-train}}
\newcommand{\rot}{\operatorname{rot}}
\newcommand{\KS}{\operatorname{KS}}
\newcommand{\ceil}[1]{\left\lceil #1\right\rceil}
\newcommand{\floor}[1]{\left\lfloor #1\right\rfloor}
\title[Diagonal Packing for Encrypted Training]{Revisiting Diagonal Packing for Non-Interactive Encrypted Neural Network Training}
\ifshowauthors
  \author{Author names to be supplied}
  \affiliation{\institution{Affiliation to be supplied}\city{City to be supplied}\country{Country to be supplied}}
  \email{author@example.invalid}
\fi
\begin{document}
\begin{abstract}
Training a neural network under homomorphic encryption requires more than an efficient forward matrix product: transposes, batch reductions, optimizer state, and repeated refreshes can dominate the cost. We revisit diagonal packing for this setting. Starting from the earlier rectangular transpose, we combine a periodic representation using the smaller padded dimension's number of diagonals with feature-major SIMD mini-batch lanes and baby-step giant-step evaluation. Compact optimizer state amortizes bootstrapping across diagonals. We implement the resulting training path in a Rust CKKS stack with instrumented CPU and CUDA arithmetic. We compare key-switch accounting conventions, bootstrap placement, and measured operator latencies, and explicitly separate layout effects from bootstrap implementations. The available encrypted subset experiments support a conditional key-switch advantage for a larger batch, but do not settle the registered accuracy and equal-security comparison. Full-setting encrypted results, GPU training ablations, and physical GPU bootstrap measurements remain pending. We report these gaps together with negative results rather than interpreting implementation completeness as research success. This interim draft establishes the algorithms and reproducible evidence needed to finalize that comparison.
\end{abstract}
\begin{CCSXML}
<ccs2012><concept><concept_id>10002978.10002979.10002980</concept_id><concept_desc>Security and privacy~Cryptography</concept_desc><concept_significance>500</concept_significance></concept><concept><concept_id>10010147.10010257.10010293.10010294</concept_id><concept_desc>Computing methodologies~Neural networks</concept_desc><concept_significance>300</concept_significance></concept></ccs2012>
\end{CCSXML}
\ccsdesc[500]{Security and privacy~Cryptography}
\ccsdesc[300]{Computing methodologies~Neural networks}
\keywords{homomorphic encryption, CKKS, encrypted training, diagonal packing, bootstrapping}
\maketitle
\input{sections/introduction}
\input{sections/preliminaries}
\input{sections/packing}
\input{sections/noninteractive}
\input{sections/gpu}
\input{sections/evaluation}
\input{sections/discussion}
\bibliographystyle{ACM-Reference-Format}
\bibliography{refs}
\appendix
\input{sections/appendix}
\end{document}
 for placeholders.
\newif\ifshowauthors
\ifdefined\AuthorDraft\showauthorstrue\else\showauthorsfalse\fi
\ifshowauthors
  \documentclass[sigconf,review]{acmart}
\else
  \documentclass[sigconf,review,anonymous]{acmart}
\fi
\setkeys{acmart.cls}{balance=false}
\setcopyright{none}
\settopmatter{printacmref=false}
\acmConference[Draft]{Unsubmitted research draft}{November 2026}{No venue assigned}
\acmDOI{}
\acmISBN{}
\usepackage{amsmath}
\usepackage{graphicx}
\usepackage{booktabs}
\newcommand{\sys}{\textsc{ckks-train}}
\newcommand{\rot}{\operatorname{rot}}
\newcommand{\KS}{\operatorname{KS}}
\newcommand{\ceil}[1]{\left\lceil #1\right\rceil}
\newcommand{\floor}[1]{\left\lfloor #1\right\rfloor}
\title[Diagonal Packing for Encrypted Training]{Revisiting Diagonal Packing for Non-Interactive Encrypted Neural Network Training}
\ifshowauthors
  \author{Author names to be supplied}
  \affiliation{\institution{Affiliation to be supplied}\city{City to be supplied}\country{Country to be supplied}}
  \email{author@example.invalid}
\fi
\begin{document}
\begin{abstract}
Training a neural network under homomorphic encryption requires more than an efficient forward matrix product: transposes, batch reductions, optimizer state, and repeated refreshes can dominate the cost. We revisit diagonal packing for this setting. Starting from the earlier rectangular transpose, we combine a periodic representation using the smaller padded dimension's number of diagonals with feature-major SIMD mini-batch lanes and baby-step giant-step evaluation. Compact optimizer state amortizes bootstrapping across diagonals. We implement the resulting training path in a Rust CKKS stack with instrumented CPU and CUDA arithmetic. We compare key-switch accounting conventions, bootstrap placement, and measured operator latencies, and explicitly separate layout effects from bootstrap implementations. The available encrypted subset experiments support a conditional key-switch advantage for a larger batch, but do not settle the registered accuracy and equal-security comparison. Full-setting encrypted results, GPU training ablations, and physical GPU bootstrap measurements remain pending. We report these gaps together with negative results rather than interpreting implementation completeness as research success. This interim draft establishes the algorithms and reproducible evidence needed to finalize that comparison.
\end{abstract}
\begin{CCSXML}
<ccs2012><concept><concept_id>10002978.10002979.10002980</concept_id><concept_desc>Security and privacy~Cryptography</concept_desc><concept_significance>500</concept_significance></concept><concept><concept_id>10010147.10010257.10010293.10010294</concept_id><concept_desc>Computing methodologies~Neural networks</concept_desc><concept_significance>300</concept_significance></concept></ccs2012>
\end{CCSXML}
\ccsdesc[500]{Security and privacy~Cryptography}
\ccsdesc[300]{Computing methodologies~Neural networks}
\keywords{homomorphic encryption, CKKS, encrypted training, diagonal packing, bootstrapping}
\maketitle
\section{Introduction}
\label{sec:intro}
Homomorphic encryption allows a service to evaluate a public computation without learning the encrypted inputs. Training is particularly demanding: the model changes at each step, gradients require transposed products, and a long optimization trajectory repeatedly exhausts the available multiplicative depth. A useful packing scheme must therefore account for the entire update, not just inference.

Mihara et al.~\cite{mihara2020} showed that diagonal-packed weights support transposition through rotations rather than ciphertext multiplications. Their interactive training design used a client to restore ciphertexts between updates. ReBoot~\cite{reboot2026} subsequently demonstrated non-interactive CKKS neural-network training, combining packed arithmetic, local learning signals, and bootstrapping. Non-interactivity itself is therefore not a novelty claim of this work. Our question is whether diagonal packing remains competitive when batch reductions, server-side layout maintenance, optimizer state, and bootstrapping are all charged.

We develop \sys{}, an MIT-licensed Rust implementation, and study three questions. \textbf{RQ1} asks how many key switches a complete training update requires compared with a ReBoot-style layout and the released ReBoot baseline. \textbf{RQ2} asks how bootstrap placement interacts with depth and compact state. \textbf{RQ3} asks what CUDA execution and its constituent optimizations contribute to measured training cost. A secondary continuity experiment revisits the original Iris setting; it is an interactive demonstration, not evidence of non-interactive training performance.

Our contributions are: (i) a rectangular, batched diagonal layout with a rotation-only transpose and explicit reduction costs; (ii) a compact-state non-interactive update whose operation counts can be checked against structural predictions; and (iii) an implementation and evaluation protocol that separates cryptographic assumptions, counter-derived comparisons, measured timing, and missing runs. Hoisting, baby-step giant-step (BSGS), sparse-secret encapsulation (SSE), and bootstrapping are established techniques; the contribution lies in their training-specific composition and accounting, not in inventing these primitives.

\paragraph{Status of the evidence.}
This is an interim full draft. The generated tables retain pending rows for the registered full-setting encrypted experiment and GPU training ablation. The completed encrypted subset runs use a different input dimension and cannot replace those rows. The released baseline and our preset also have different secret distributions and security estimates. We consequently do not claim that the registered equal-security, accuracy-matched research goal has been met. Section~\ref{sec:limits} identifies what would change that conclusion.

\section{Preliminaries and Threat Model}
\label{sec:prelim}
\paragraph{Arithmetic.}
CKKS~\cite{ckks2017} encodes approximate vectors in
$R_Q=\mathbb{Z}_Q[X]/(X^N+1)$, with up to $N/2$ complex slots. We use real slot values. A ciphertext at level $\ell$ uses $q_0,\ldots,q_\ell$ and a scale $\Delta$. Multiplication followed by rescaling consumes a modulus prime. Rotations permute slots through a Galois automorphism and a key switch. Hybrid RNS key switching uses special primes $P$ and a digit decomposition; its cost depends on the active limbs, digit sizes, and memory traffic.

We count rotations (including conjugations), relinearizations, and explicit SSE switches:
\begin{equation}
  \KS = n_{\mathrm{rot}}+n_{\mathrm{relin}}+n_{\mathrm{SSE}}.
  \label{eq:ks}
\end{equation}
A ciphertext tensor product is not automatically a key switch: lazy relinearization may combine several products before one switch. Bootstrap events are recorded separately and their constituent operations are already in Eq.~\eqref{eq:ks}; adding a bootstrap-equivalent charge again would double count. Reported counts cover the whole network and are averaged over processed examples, not over layers. A count is a workload metric, not a machine-independent latency estimate.

\paragraph{Refreshes.}
An interactive refresh decrypts and re-encrypts at the key owner. A bootstrap instead evaluates a refresh circuit using public evaluation keys. Our non-interactive path uses SSE~\cite{sse2022}, centered modulus raising, coefficient-to-slot and slot-to-coefficient transforms, and an approximate modular-reduction polynomial. Public transforms use collapsed stages and BSGS; exact scale tracking makes them composable with training.

\paragraph{Protected objects and public leakage.}
We consider a semi-honest evaluation server and a single trusted key owner, under the usual RLWE and evaluation-key assumptions of the chosen CKKS construction. No new security reduction is claimed. Training examples, labels, model weights, intermediate activations, gradients, and optimizer state are encrypted. The architecture, dimensions, batch schedule, hyperparameters, parameters, evaluation keys, and algorithmic control flow are public. The key owner initializes and encrypts the data and initial model, and decrypts the resulting model. Known initialization distributions are not themselves secret. This model does not hide dimensions, resource use, access schedules, or timing, and it does not provide malicious-server verification, differential privacy, multi-party key custody, or protection against inference from a legitimately decrypted model.

\paragraph{Diagnostics are not a server capability.}
The benchmark process can also instantiate a secret-key holder, decrypt checkpoints and predictions, compare against a plaintext twin, and report accuracy or activation ranges. Those are client-side diagnostics for a trusted experiment harness; they must not be exposed as server logs in a deployment. The generic non-interactive trainer has no secret key and makes public depth-based refresh decisions. Disabling diagnostic decryption changes neither the encrypted arithmetic nor the public schedule. The interactive path explicitly involves the client during training and is labeled separately. The in-library ReBoot-layout microbenchmark uses plaintext labels and client-side optimizer updates: it supports arithmetic/count comparisons, not the full confidentiality claim or a faithful non-interactive baseline.

\section{Rotation-Only Transposition for Mini-Batch Training}
\label{sec:packing}
\subsection{A common layout for activations and weights}
Let $B$ be a power-of-two lane count and $d$ a power-of-two feature capacity such that $dB$ divides the slot count. Slot $jB+b$ stores feature $j$ of sample $b$. Blocks repeat when $dB$ is smaller than the slot count. A layer has positive integer input and output dimensions $d_i,d_o$. Define the padded dimensions $n_i=2^{\ceil{\log_2 d_i}}$ and $n_o=2^{\ceil{\log_2 d_o}}$, and require $n_i,n_o\le d$. A vector of padded dimension $n\in\{n_i,n_o\}$ is canonical when its feature sequence is $n$-periodic in every lane; unused logical features are zero. This periodicity eliminates a separate replication step at tall layers.

Zero-pad the original $d_o\times d_i$ weight matrix to $n_o\times n_i$. Let $m=\min(n_o,n_i)$ and $L=\max(n_o,n_i)$. Both divide $d$, so the $L$-periodic diagonals, rotations modulo $LB$, and integer logarithmic folds below are well-defined. Store the generalized diagonals
\begin{equation}
 D_c[j]=W[j\bmod n_o,(j+c)\bmod n_i],\quad 0\le c<m.
 \label{eq:diagonal}
\end{equation}
Weights are shared across all lanes. With BSGS split $m=gu$ and $c=ga+b$, the stored diagonal is
$P_c=\rot(D_c,-gaB)$. The encrypted product is
\begin{equation}
 y=\sum_{a<u}\rot\!\left(\sum_{b<g}P_{ga+b}\odot\rot(x,bB),gaB\right).
 \label{eq:bsgs}
\end{equation}
Here a positive rotation reads slot $s+k$ into output slot $s$; $\odot$ denotes a slot-wise product. For a wide layer, $\log_2(n_i/n_o)$ rotations and additions fold partial outputs together. For a tall layer, the canonical input already has the required replication. Neither case adds a mask to the fold.

\subsection{Transpose without products or masks}
Rectangular diagonal transposition without ciphertext multiplications is inherited from Mihara et al.~\cite{mihara2020}. Their general method uses rotations and additions; their implementation also reports OpenMP batch processing. Here we compose that idea with the periodic $m$-diagonal representation, feature-major SIMD lanes, and BSGS pre-rotations.

The transpose diagonal $E_i[j]=W[(j+i)\bmod n_o,j\bmod n_i]$ equals a rotation of one $D_c$. Choose $c=(-i)\bmod m$, and set $t=-c$ for a wide layer or $t=i$ for a tall layer. In stored BSGS form,
\begin{equation}
 Q_i=\rot(P_c,\,[t-g\floor{i/g}+g\floor{c/g}]B).
 \label{eq:transpose}
\end{equation}
The rotation amount is taken modulo $LB$. Thus transposition uses neither a ciphertext product nor a mask and consumes no multiplicative level. It is \emph{not} free: each nonzero rotation requires a key switch. Dedicated keys minimize rotation count; composing existing baby/giant keys saves key memory at the expense of extra switches. Our training runs use key reuse, and the extra work is included in the counters.

\subsection{Gradients, reductions, and a symbolic cost model}
For layer gradient $\delta$ and input $a$, a pre-rotated weight-gradient diagonal is
\begin{equation}
 \nabla P_{ga+b}
 =\operatorname{batchsum}\bigl(\rot(\delta,-gaB)\odot\rot(a,bB)\bigr).
 \label{eq:gradient}
\end{equation}
Batch summation first reduces lanes by powers-of-two rotations. A lane-zero mask isolates the reduced value; inverse-direction rotations replicate it when replicated gradients are required. The mask has nonintegral polynomial coefficients in the CKKS embedding, so encoding it at unit scale is not an exact, zero-depth shortcut. It consumes a level. Constants for the learning rate and batch normalization are folded into a necessary plaintext multiplication at its exact target scale.

\begin{table}[t]
\centering\small
\caption{Symbolic layer costs for the implemented layout, before any wide-layer fold. $g,u$ are the BSGS split, $m=gu$. Gradient costs show the replicated interactive form; compact non-interactive state changes the final replication.}
\label{tab:symbolic}
\begin{tabular}{lcc}
\toprule
Operation & Rotations & Relins \\
\midrule
BSGS encrypted matvec & $g+u-2$ & $u$ \\
Dedicated transpose & $m-1$ & $0$ \\
Replicated weight gradient & $g+u-2+2m\log_2 B$ & $m$ \\
\bottomrule
\end{tabular}
\end{table}
Hoisting reuses decomposition for rotations of the same input~\cite{halevi2018}. It changes wall-clock work but not Eq.~\eqref{eq:ks}. Lazy relinearization applies within a matvec giant-step sum. It cannot remove the separate relinearizations needed before gradient-diagonal rotations. These distinctions matter when interpreting ablation results.

\paragraph{Extensions relative to the interactive design.}
The first two extensions studied here are the periodic $\min(n_o,n_i)$-diagonal representation with feature-major SIMD mini-batch lanes, and its composition with hoisting and BSGS. General rectangular transposition and batch processing themselves are inherited capabilities, not novelty claims. The remaining four extensions---bootstrap placement, polynomial activations, momentum, and loss conventions---are specified in Section~\ref{sec:noninteractive}. None removes the cost of maintaining canonical server-side state; the initial interactive cost model omitted some replication work and was therefore too optimistic.

\section{Non-Interactive Training}
\label{sec:noninteractive}
\subsection{Compact state and the complete update}
Refreshing every replicated diagonal separately wastes bootstrap slots. We instead reduce each gradient to one lane per feature and combine masking, scaling, and lane shifting in a plaintext product. Selected lanes from several diagonals form compact weight and momentum containers. A public compact-layout plan fixes their offsets and the required rotations; unpacking restores the periodic diagonals for the next forward pass. Bias slots can share unused compact positions. Masked packing and unpacking both have a depth cost and are included in the closed-form predictor.

The server performs encrypted forward propagation, forms the loss gradient from encrypted labels, propagates through the transposed weights, and accumulates gradients over micro-batches before one optimizer update. Increasing the effective batch beyond the physical lane count therefore amortizes optimizer refreshes without pretending the examples fit in one ciphertext. It also changes optimization behavior and cannot be treated as a cost-only substitution.

The original dense-secret bootstrap constants did not satisfy the desired precision and depth budget. The implemented preset instead uses a dense ternary main key and a temporary sparse key during modulus raising. The SSE pair is returned to the dense main key before the remaining arithmetic. Collapsed transforms and a level-optimal Chebyshev evaluator fit the selected modulus bands. This is an engineering parameter choice, not a new bootstrap construction or an inherited failure-probability theorem. Appendix~\ref{app:params} specifies the chain and the exposed key instances.

\subsection{Three refresh schedules (extension 3)}
\texttt{AfterUpdate} (AU) bootstraps every compact weight and velocity container after each update. \texttt{WhenNeeded} (WN), an additional schedule beyond the original plan, does so only if the next update would exceed the remaining depth. The decision uses the next batch's fullness: padded lanes at a linear output need a validity mask, so a partial batch can consume an extra level. An unknown next shape is conservatively treated as partial.

\texttt{AfterActivation} (AA) refreshes each hidden activation and retains the WN rule for compact weights. For the shallow eMLP-1 workload, activation refresh does not reduce the frequency of weight refresh beyond WN; it adds a bootstrap per hidden activation and micro-batch. Its motivation is enabling deeper circuits that otherwise exceed the budget, not an automatic speedup. Activation magnitudes must fit the supported bootstrap range. The observed AA subset run exceeded the conservative range in some slots and had a larger activation error; it was not clipped, and does not certify arbitrary activation distributions.

For a sequence of identical updates consuming $D$ levels with $L_a$ usable levels after unpacking, $\floor{L_a/D}$ is only an initial scheduling bound. Unpack costs between updates, partial batches, and the actual state level must also enter the scheduler. We check the structural predictor against every measured update rather than replacing this stateful schedule with a fixed average.

\subsection{Activation, optimizer, and loss (extensions 4--6)}
\paragraph{Polynomial activation.}
We retain the earlier squared activation for continuity and implement $z^2+z$ for the ReBoot comparison. A polynomial quartic option is available for depth experiments. Its additional depth and available tests do not establish a measured accuracy improvement; that planned claim remains unsupported.

\paragraph{Momentum.}
Both SGD and Nesterov-style momentum are supported. Writing $h=\eta(1+\mu)G/n$ and scaled velocity $U=\eta(1+\mu)V$ gives
\begin{align}
 U'&=\mu U+h,\nonumber\\
 W'&=W-h-\frac{\mu^2}{1+\mu}U.
 \label{eq:nag}
\end{align}
Constants are folded into exact-scale products already needed by the gradient path. A changing learning rate requires rescaling the stored velocity representation, and its lane shifts are charged. Biases use SGD; the faithful ReBoot configuration freezes their encrypted zero values while retaining the implementation's bias arithmetic. Weight decay is not implemented; the captured faithful setting uses zero decay.

\paragraph{Loss.}
A squared-error loss yields an affine output gradient; MSE normalization and ReBoot's sum reduction are matched through the learning-rate convention. The faithful configuration uses the released initialization, sample order, output padding, and schedule, not just the paper's nominal hyperparameter names. Softmax/cross-entropy approximation is optional in the plan and is not implemented or evaluated here. Block MLPs follow local classifier losses and stop gradients across blocks; they must not be described as an ordinary deeper back-propagation network.

\section{GPU Implementation}
\label{sec:gpu}
\sys{} separates RNS arithmetic, CKKS evaluation, neural-network operations, CUDA execution, and benchmark orchestration. Generic \texttt{HeOps} and \texttt{BootstrapOps} interfaces let the same trainer use CPU or device ciphertexts. The evaluator checks parameter identity, levels, scales, key availability, and device ownership before work. Counters commit after success; the CUDA bootstrap also waits for final device synchronization before merging its internal counters.

The CUDA backend uses device NTTs, basis conversion, pointwise arithmetic, hybrid key switching, and BSGS matvecs. Small public launch data use pinned transfer buffers. Reusing a hoisted decomposition avoids repeating the expensive preparation for related rotations. Device/CPU bit parity is meaningful only for cases actually executed; memory-skipped test functions can still report process success.

\paragraph{Streams and memory.}
Independent products may be issued to several streams, but a training step's dependent matvecs are executed sequentially. The measured pre-fix multi-stream benchmark was slower than sequential execution (Table~\ref{tab:gpu_vs_cpu_ops}). Calling-thread submission and pool handling were subsequently revised; their performance remeasurement is pending. An effective-bandwidth estimate for addition does not prove matvec occupancy or full-device saturation. That explanation remains a hypothesis until profiled.

The shared device has limited available memory. Training can upload keys capped to an explicitly selected level; this preserves the computation at that level, but does not measure the top-level workload. Allocation preflight uses estimated resident/key/workspace sizes. Pool trims occur after large setup stages, at epoch boundaries, and between ablation configurations. Estimated need and pool checkpoints are not a measured memory peak, and a successful preflight does not guarantee later allocations under contention.

\paragraph{Bootstrap boundary.}
The CUDA bootstrap implementation keeps ciphertext arithmetic on the device through SSE, modulus raising, transforms, and EvalMod. A narrow shared polynomial executor preserves the CPU operation order and rounding. Public integer constants are still materialized on the host and uploaded; future device timing must include that overhead. Local CPU replay and CUDA compilation support code readiness, not physical-GPU precision, latency, or memory feasibility. Those hardware measurements are pending.

The training benchmark still uses a CPU bootstrap adapter. Level operations run on the GPU; non-interactive refresh downloads to the CPU and uploads the result. A compiled device bootstrap API has not changed that benchmark. Its timing must therefore separate device work, transfers, and CPU refresh. Likewise interactive GPU training uses a client refresh. Neither path justifies an end-to-end all-GPU training claim.

\section{Evaluation}
\label{sec:evaluation}
\subsection{Protocol and evidence classes}
We distinguish a released baseline execution, an in-library layout microbenchmark, and a cost model. The first supplies observed wall-clock and bootstrap events; the second isolates arithmetic layout counts under our evaluator; the third supplies uninstrumented baseline quantities. A counter-derived normalization is not a separate execution. Historical generated tables retain the working name \emph{hekit} for \sys{} and its crate names.

The headline uses MNIST reduced to $196$ inputs and the released full training setting. FULL-1 has completed at $B=64$ with one initial-weight seed; FULL-2 at the released baseline batch remains pending. FULL-1 is partial full-setting evidence, not a substitute for the FULL-2 accuracy gate. The completed faithful encrypted subset uses $784$ inputs, the baseline's sample indices and initial weights, and its training conventions. Output width is padded to $16$, including the unused classes in the released loss; replacing this with a width-$10$ loss is not assumed equivalent. Earlier ``legacy'' runs use a different learning rate, bias policy, and initialization. Their counters and timing remain useful, but their accuracy is not comparable to the released baseline. The last subset batch is partial, affecting depth, refresh scheduling, and the effective sum-versus-mean learning-rate conversion.

The paired server measurements use the same CPU and thread cap, with different CPU parallel runtimes. Jobs are not core-pinned, and some rows overlap unrelated shared-host load. Client work and setup scopes also differ: baseline step time includes batch encryption and prediction decryption, whereas our update time excludes those separately reported client operations; both exclude key setup. Thus even a same-host ratio is not a controlled architectural speedup. We do not pool CPU and GPU samples. Encrypted rows have a single seed unless the table states otherwise; within-run iterations and Criterion samples are not independent seeds.

\subsection{RQ1: key switches and wall-clock}
\input{tables/rq1_ks}
\input{tables/wallclock_same_cpu}
Let $K$ be measured total key switches, $H$ bootstrap events, and $S$ examples processed. Let $c_h$ be our measured per-bootstrap cost and $c_r$ the baseline bootstrap model; let $R$ denote the baseline non-bootstrap count per example, $r$ its refreshes per update, and $B_r$ its own batch. The four reported ratios are
\begin{align}
 \rho_a&=\frac{K/S}{R+rc_r/B_r},&
 \rho_b&=\frac{(K-c_hH+c_rH)/S}{R+rc_r/B_r},\nonumber\\
 \rho_c&=\frac{K/S}{R+rc_h/B_r},&
 \rho_d&=\frac{(K-c_hH)/S}{R}.
 \label{eq:ratios}
\end{align}
Convention (a) uses each side's own bootstrap cost. Convention (b) charges both the baseline cost; (c) charges both our cost; (d) excludes bootstrap arithmetic on both sides. The interactive comparison is separate from (d): omitting bootstrap arithmetic does not remove non-interactive compact packing and unpacking. ReBoot remains at its own batch in the denominator; a hypothetical same-batch denominator would answer a different question.

The baseline's bootstrap cost remains a model because the measured OpenFHE execution does not expose key-switch counters. Its observed bootstrap frequency is used, not replaced with our schedule. The subset non-bootstrap reference is the in-library Lazy layout measurement; the full-setting reference is still a model. These boundaries are visible in Table~\ref{tab:rq1_ks}.

\input{figs/rq1_ratios}
The faithful larger-batch subset row meets the registered count threshold under every reported convention, while the smaller batch meets it only under the own-bootstrap convention. The generated figure makes that dependence visible. The corresponding larger-batch timing is also favorable on the shared CPU (Table~\ref{tab:wallclock_same_cpu}), but its subset accuracy drops relative to the baseline, including in the plaintext twin. Consequently this observation is a tradeoff, not a successful accuracy-matched RQ1 result. FULL-1 also meets the count threshold under all four conventions. Its completed training time has no measured full-setting ReBoot CPU denominator. The pending FULL-2 accuracy gate, missing baseline timing, and security mismatch below prevent a final verdict.

\subsection{RQ2: placement and depth}
The AU, WN, and AA subset rows in Tables~\ref{tab:rq1_ks} and~\ref{tab:wallclock_same_cpu} provide the placement comparison. For matched legacy eMLP-1 settings, WN reduces bootstrap events and total time relative to AU. AA leaves the same lazy weight-refresh frequency and adds activation refreshes, so it does not improve this shallow workload. Its benefit is a feasible depth budget for some deeper polynomial circuits, checked separately as a structural test rather than an observed training-accuracy gain.

The corrected structural predictor equals the measured counters at every completed update. This validates accounting for the executed paths; it does not validate the baseline bootstrap model or a convergence claim. The initial cost model missed replication, key-reuse, bias, and schedule-related terms, and its optimistic interactive estimate is superseded by that explicit accounting. FULL-1 extends exact schedule agreement to its complete registered update sequence and records accumulated weight and bias error against its same-batch plaintext twin. That observed error is larger than in the earlier subset checks; close final accuracy does not imply a uniform small parameter error. FULL-2 and guarantees across seeds or arbitrary trajectories remain unmeasured.

\subsection{Accuracy, local blocks, and security}
\input{tables/accuracy_vs_reboot}
\input{tables/plain_datasets}
\input{tables/security}
Table~\ref{tab:accuracy_vs_reboot} separates the released full-setting reference, subset results, and published paper accuracy. The released defaults and paper-setting reconstruction are different protocols. We do not choose whichever reference makes the result look best. The registered tolerance is tested by FULL-2 against the captured released full-setting seed-zero reference; that encrypted row has not completed. FULL-1 reports its larger-batch accuracy and same-batch twin separately, without substituting either for this gate. Matching a plaintext twin is necessary to diagnose arithmetic fidelity, but is insufficient to establish baseline-equivalent generalization.

Plaintext dataset experiments (Table~\ref{tab:plain_datasets}) provide broader optimization checks. MNIST and Fashion-MNIST reproduce the stated comparison tolerance at the fixed final epoch. Breast Cancer does not: the fixed-epoch min-max result misses its target, and selecting the best test epoch is optimistic leakage from model selection. The z-score setting diverges for the polynomial activation. Fashion-MNIST and Breast Cancer results are plaintext twins, not encrypted training measurements. The local-block plaintext twins reproduce the released eMLP-2/eMLP-3 values for the recorded seed; their encrypted subset rows are pending. The block experiments do not extend the current headline claim to deeper networks.

Table~\ref{tab:security} applies common estimator models to the actual exposed key moduli and secret distributions. The evaluated main and SSE instances of our preset exceed the chosen operation-count target under those models; the released ReBoot preset does not. This is an estimate over the documented attack search, not a security proof. The smallest sparse-secret values use attacks with extreme memory requirements, and even non-MITM minima have substantial or unreported memory demands. These are not practical attack demonstrations. Equal $N$ and a requested library security label do not establish equal security. Results must remain conditional until either a security-different comparison is explicitly accepted or an independently estimated and remeasured alternative baseline is available.

\subsection{RQ3: device operations and ablation}
\input{tables/gpu_vs_cpu_ops}
\input{tables/rq3_ablation}
Table~\ref{tab:gpu_vs_cpu_ops} contains synchronized operation latencies, including capped-level rows and memory skips. The HEonGPU relinearization reference~\cite{gpuKeys2025} is nominal: its modulus count differs. A linear-in-limb adjustment is a cost-model extrapolation, not an equal-parameter measurement. The published table value is inherited from the archived operator report; the bibliographic metadata was checked for this draft, but a fresh PDF fetch was blocked. The later pool/submission fix has not been retimed.

Table~\ref{tab:rq3_ablation} deliberately leaves the server rows pending. The insecure small-parameter CPU smoke validates counts and checks only; its high-load timings are inadmissible for performance comparisons. The flag named \texttt{bsgs=false} keeps pre-rotated weights rather than implementing the strongest textbook diagonal baseline. Thus its BSGS gain would be an upper bound against that alternative, while the marginal hoisting and lazy-relinearization gains within those rows would be lower bounds. Hoisting changes no operation counter; the gradient outer product retains layout-required rotations independent of the BSGS flag. All four caveats travel with the table.

CPU and GPU training comparisons require the same optimizer, level, layout, and batch. Registered SGD fallback rows are compared only with the corresponding CPU SGD row. A resource skip is unmeasured; a device OOM is a failure. Test-process success cannot stand in for an actually executed production-size parity gate. GPU bootstrap precision, latency, allocation behavior, and peak memory are still unmeasured. HEonGPU's separate bootstrap reference~\cite{heongpuBench} has a different remaining-level budget and main-secret distribution; the fixed plan target is not a normalized speed-ratio result. RQ3's epoch-time question is therefore unanswered.

\section{Related Work, Limitations, and Conclusion}
\label{sec:limits}
\subsection{Positioning}
The earlier diagonal-training design~\cite{mihara2020} motivates rotation-only transposition. We generalize the layout and charge server-side batch maintenance, compact refresh, and optimizer state. ReBoot~\cite{reboot2026} already establishes non-interactive training and local learning; its released implementation is our application baseline, not a source of imported implementation code. HETAL~\cite{hetal2023} studies encrypted transfer learning and efficient matrix products in a different training scope. Our evaluation does not reproduce its workloads or establish a ranking against it.

The cryptographic building blocks are efficient linear transforms and hoisting~\cite{halevi2018}, SSE~\cite{sse2022}, and CKKS~\cite{ckks2017}. SEAL~\cite{seal} provides interoperability checks; OpenFHE~\cite{openfhe2022} supplies the separately executed baseline. HEonGPU~\cite{gpuKeys2025,heongpuBench} is a device-operation reference, not an end-to-end training competitor measured under our protocol. An MIT implementation improves reuse, but licensing and pipeline completeness are not evidence of a cryptographic or learning-performance advance. The approved dependency policy permits MIT, Apache, BSD, and ISC licenses, with one narrow Unicode-3.0 character-table exception for \texttt{unicode-ident} version \texttt{1.0.27}. That component requires its complete Unicode notice; this is not a general Unicode-license exception.

\subsection{Limits and missed goals}
\paragraph{Scope and comparison.}
FULL-1 and FULL-2 supply completed full-setting encrypted runs at two batches with the same initial-weight seed, including their schedules and accumulated error. FULL-2 passes the registered baseline-batch accuracy gate, but misses the common-baseline-bootstrap count target. FULL-1's successful count comparisons are a separate configuration. Full-setting ReBoot CPU wall-clock is unmeasured, and the common-security condition is unmet. Neither these two configurations nor the earlier subsets establish multi-seed encrypted behavior; multi-seed plaintext results are a different evidence class. The observed parameter error does not establish a general error bound. Batch choices trade refresh amortization against learning dynamics. The effective baseline security differs, initialization and bias arithmetic require care, and legacy rows cannot support accuracy claims. ReBoot's in-library layout timing was measured on a different machine and is not a same-CPU baseline; its public-label/client-update path has a weaker confidentiality contract. The official encrypted block-model baselines have not been run: their reported memory demand exceeds the available host budget.

\paragraph{Cost and correctness.}
A key switch at a low level need not cost the same as one at the top level; counts omit memory traffic, cache effects, and key residency. The initial compact-layout assumptions required actual CKKS tests and the first model was optimistic. Replicated copies were once unstable during interactive training; canonical-layout refresh was required to avoid divergence. The earlier paper's published timing total and phase sum disagree, so the continuity demonstration does not silently reconcile them or infer a server-only speedup. The original CPU NTT target and several SEAL-relative key-switch performance targets were missed. The early CUDA NTT spike was not timed against its registered hardware reference; later device-operator measurements do not retroactively supply that measurement. Some SEAL comparisons also have different special-prime/key-switch decompositions. We report these as missed performance targets, not failed numerical gates.

\paragraph{Numerical limits.}
SSE and modulus-band changes were needed because the planned dense-secret overflow range and original prime allocation did not meet the bootstrap precision budget. Unit-scale nonconstant masks were invalid, and exact-scale folding is required to fit the training circuit. The accepted security claim uses specified sieve cost models, not a universal security level; Core-SVP estimates are lower. Dense-secret hybrid search is bounded rather than an exhaustive global optimization. An all-sparse main-key alternative did not meet the selected security criterion and is not used. KAT coverage is limited to the supported SEAL preset, with asymmetric comparison directions; it does not independently validate our bootstrap or larger hybrid preset. Standalone OpenFHE/Lattigo bootstrap precision parity remains absent. Random full-slot CPU precision checks do not prove a negligible failure probability for arbitrary training distributions or activations outside range.

\paragraph{Learning limits.}
The quartic activation's claimed accuracy benefit is unmeasured. The evaluated encrypted trainer omits weight decay and optional softmax/cross-entropy. The separate plaintext paper-setting reconstruction includes weight decay; it is not the faithful encrypted protocol. Bias momentum is not matched to NAG, and faithful bias freezing preserves extra encrypted arithmetic. Fixed-final-epoch Breast Cancer accuracy misses its target; test-selected best epochs are not a remedy. No encrypted Fashion-MNIST or Breast Cancer accuracy is reported. Higher-depth local-block training has plaintext fidelity checks but no completed encrypted comparison. The continuity Iris experiment is interactive and differs in padding, randomization, and normalization; it cannot establish the primary non-interactive claim. Threshold or multi-party FHE, scheme switching, transformer training, and Metal execution are outside the evaluated scope.

\paragraph{Device and implementation limits.}
Full-key GPU training, the matched training parity gate, the RQ3 timing ablation, and GPU bootstrap measurements are pending. Repaired multi-stream latency is measured only for independent products at a capped level; it cannot supply those missing training results. Memory estimates and pool snapshots do not establish a peak. All-off flags do not implement the strongest naive diagonal algorithm. The current training benchmark still uses host bootstrap fallback. Device error handling and numerical replay are useful engineering evidence, but compile-only paths and skipped tests are not physical hardware results. CPU/GPU side-channel resistance, constant-time sampling, secret zeroization, malformed-input hardening, and independent cryptographic audit are outside the current implementation's assurance. Some internal invariants are debug assertions; decode/parameter mismatches can panic. The artifact is research software, not a deployment security claim.

\subsection{Conclusion}
Batched diagonal packing preserves a rotation-only transpose while supporting compact non-interactive optimizer updates. The completed full-setting runs show a larger-batch count advantage and baseline-batch accuracy within tolerance, together with a missed normalized count target at the baseline batch. The security mismatch and unmeasured full-setting baseline time prevent an unconditional superiority claim. The device microbenchmark also shows that stream and pool choices matter, without resolving training performance. The resulting evidence supports a conditional re-evaluation of the layout and identifies the matched-security, baseline-timing, and hardware experiments still needed.

\bibliographystyle{ACM-Reference-Format}
\bibliography{refs}
\appendix
\section{Parameters, Security, and Reproduction}
\label{app:params}
\subsection{Implemented chain and bootstrap contract}
The production bootstrap preset uses $N=2^{16}$, $N/2$ real slots, scale $2^{42}$, and Q-prime bit sizes
\begin{equation*}
 [54,\underbrace{42,\ldots,42}_{16},\underbrace{34,\ldots,34}_{4},
 \underbrace{58,\ldots,58}_{7},\underbrace{56,\ldots,56}_{4}].
\end{equation*}
Four special primes have bit size $60$, with eight hybrid digits. The bands allocate four levels to each linear transform and seven to EvalMod, leaving level $16$ before training unpack. EvalMod uses Chebyshev degree $127$, interval parameter $K=13$, and no double-angle step. These are implemented design constants; the superseded original-plan bands and dense-secret range are not used.

The bootstrap requires level zero and exactly the preset scale before work. Dense-to-sparse SSE is restricted to $q_0P$ under a temporary key of weight $32$; the return key and remaining arithmetic use the dense main key over divisors of $QP$. CPU full-slot tests exceed the registered precision threshold for random real inputs in $[-1,1]$. This does not imply a worst-case failure bound. The GPU implementation has local replay and compilation coverage, but its physical precision and memory feasibility remain unmeasured.

\subsection{What the estimator models}
Our default sampler gives $-1,0,1$ probabilities $1/4,1/2,1/4$. We model the dense secret with a fixed-count sparse-ternary proxy rather than a uniform-ternary distribution; SSE is a separate small-modulus instance. The reported attack input rounds the exposed $\log_2(QP)$ upward and uses $m=N$ samples. Dense and sparse main keys must not be interchanged at the same modulus without re-estimation.

The estimator is pinned to commit \texttt{ae45d740ceb1}; methodology follows Albrecht et al.~\cite{albrecht2015} and the maintained estimator documentation~\cite{latticeEstimator}. The recorded search includes primal/dual and selected hybrid/MITM variants. Dense primal-hybrid optimization did not finish globally; fixed guessing-dimension scans supplement it. Arora--GB and BKW were not included. For the released ReBoot context, an observed sign split from one generated key is an input instance, not a distributional guarantee for all keys. The minima are conditional on the selected search, cost model, and memory assumptions. Table~\ref{tab:security} is not a proof, nor evidence of a practical compromise.

\subsection{Transpose identity and reduction boundary}
For Eq.~\eqref{eq:transpose}, matching row and column indices gives $t\equiv i\pmod {n_o}$ and $t+c\equiv0\pmod {n_i}$. The wide/tall cases in Section~\ref{sec:packing} solve both conditions. Substituting the pre-rotations on the source and destination yields the displayed offset. Because every weight diagonal is periodic with period $L$ in feature coordinates and all offsets are multiples of $B$, samples are not mixed. The identity therefore applies to rectangular padded layers and repeated blocks as well as square matrices. A zero offset is a copy; other offsets require the selected dedicated or composed rotation keys.

Lane reduction is different: offsets not divisible by $B$ mix lanes intentionally, and isolating one lane requires a mask. Ignoring that mask or its replication work changes both the depth and count model. A partial-batch validity mask must also exclude unused lanes from the loss derivative; zero input alone is insufficient when the model has a bias or a linear output loss.

\subsection{KAT and artifact boundaries}
The SEAL compatibility checks compare encoding, serialized ciphertext interoperability, decryption, and supported evaluation paths under the shared preset. Some directions assert exact residues while others only assert a decoded-error bound. Rotation arithmetic intentionally differs in its centered lifting. Neither this KAT nor CPU/GPU bit equality establishes security; the latter only checks implementation agreement for executed cases. Small insecure presets are smoke tests and never enter the security claim.

Tables are generated by \texttt{tools/paper\_tables.py}; field-level provenance is available through its \texttt{--sources} option. The figure generator consumes those validated table cells and writes a source manifest containing paths, fields, and SHA-256 digests. It does not rewrite result JSON. Missing registered runs produce explicit pending rows; existing invalid inputs must fail validation. Timed runs, setup, client diagnostics, and model-only extrapolations stay distinguishable. Running \texttt{scripts/run\_all.sh --tables-only} regenerates the retained tables. This does not execute the complete registered experiment matrix, whose end-to-end runner remains unfinished; a quick pipeline check is not a source of research results.

The anonymous review mode omits author metadata and identifying artifact links. ReBoot and OpenFHE are executed as external baselines; no implementation source is copied into this artifact. The estimator and TeX compiler are isolated development tools, not library dependencies. The library's dependency policy, including the specific character-table exception disclosed in Section~\ref{sec:limits}, is separate from the licenses of those external research tools. Instructions for rebuilding anonymous review, named review, and named preprint variants are provided with the source. Authorship is confirmed for this preprint. Before arXiv submission, the final source, metadata, and server-generated PDF require checking against the retained evidence; unavailable measurements remain unavailable. A later conference submission requires a separate anonymity and venue-rules review; the WAHC~2026 limit is only the provisional formatting constraint~\cite{wahc2026}.

\ifdefined\ArxivDraft
Code and evidence: \url{https://github.com/kenmaro3/ckks-train}. Release \texttt{v1.0.0-preview.1} is an earlier snapshot; this preprint has its own source revision.
\fi

\end{document}
 for placeholders.
\newif\ifshowauthors
\ifdefined\AuthorDraft\showauthorstrue\else\showauthorsfalse\fi
\ifshowauthors
  \documentclass[sigconf,review]{acmart}
\else
  \documentclass[sigconf,review,anonymous]{acmart}
\fi
\setkeys{acmart.cls}{balance=false}
\setcopyright{none}
\settopmatter{printacmref=false}
\acmConference[Draft]{Unsubmitted research draft}{November 2026}{No venue assigned}
\acmDOI{}
\acmISBN{}
\usepackage{amsmath}
\usepackage{graphicx}
\usepackage{booktabs}
\newcommand{\sys}{\textsc{ckks-train}}
\newcommand{\rot}{\operatorname{rot}}
\newcommand{\KS}{\operatorname{KS}}
\newcommand{\ceil}[1]{\left\lceil #1\right\rceil}
\newcommand{\floor}[1]{\left\lfloor #1\right\rfloor}
\title[Diagonal Packing for Encrypted Training]{Revisiting Diagonal Packing for Non-Interactive Encrypted Neural Network Training}
\ifshowauthors
  \author{Author names to be supplied}
  \affiliation{\institution{Affiliation to be supplied}\city{City to be supplied}\country{Country to be supplied}}
  \email{author@example.invalid}
\fi
\begin{document}
\begin{abstract}
Training a neural network under homomorphic encryption requires more than an efficient forward matrix product: transposes, batch reductions, optimizer state, and repeated refreshes can dominate the cost. We revisit diagonal packing for this setting. Starting from the earlier rectangular transpose, we combine a periodic representation using the smaller padded dimension's number of diagonals with feature-major SIMD mini-batch lanes and baby-step giant-step evaluation. Compact optimizer state amortizes bootstrapping across diagonals. We implement the resulting training path in a Rust CKKS stack with instrumented CPU and CUDA arithmetic. We compare key-switch accounting conventions, bootstrap placement, and measured operator latencies, and explicitly separate layout effects from bootstrap implementations. The available encrypted subset experiments support a conditional key-switch advantage for a larger batch, but do not settle the registered accuracy and equal-security comparison. Full-setting encrypted results, GPU training ablations, and physical GPU bootstrap measurements remain pending. We report these gaps together with negative results rather than interpreting implementation completeness as research success. This interim draft establishes the algorithms and reproducible evidence needed to finalize that comparison.
\end{abstract}
\begin{CCSXML}
<ccs2012><concept><concept_id>10002978.10002979.10002980</concept_id><concept_desc>Security and privacy~Cryptography</concept_desc><concept_significance>500</concept_significance></concept><concept><concept_id>10010147.10010257.10010293.10010294</concept_id><concept_desc>Computing methodologies~Neural networks</concept_desc><concept_significance>300</concept_significance></concept></ccs2012>
\end{CCSXML}
\ccsdesc[500]{Security and privacy~Cryptography}
\ccsdesc[300]{Computing methodologies~Neural networks}
\keywords{homomorphic encryption, CKKS, encrypted training, diagonal packing, bootstrapping}
\maketitle
\section{Introduction}
\label{sec:intro}
Homomorphic encryption allows a service to evaluate a public computation without learning the encrypted inputs. Training is particularly demanding: the model changes at each step, gradients require transposed products, and a long optimization trajectory repeatedly exhausts the available multiplicative depth. A useful packing scheme must therefore account for the entire update, not just inference.

Mihara et al.~\cite{mihara2020} showed that diagonal-packed weights support transposition through rotations rather than ciphertext multiplications. Their interactive training design used a client to restore ciphertexts between updates. ReBoot~\cite{reboot2026} subsequently demonstrated non-interactive CKKS neural-network training, combining packed arithmetic, local learning signals, and bootstrapping. Non-interactivity itself is therefore not a novelty claim of this work. Our question is whether diagonal packing remains competitive when batch reductions, server-side layout maintenance, optimizer state, and bootstrapping are all charged.

We develop \sys{}, an MIT-licensed Rust implementation, and study three questions. \textbf{RQ1} asks how many key switches a complete training update requires compared with a ReBoot-style layout and the released ReBoot baseline. \textbf{RQ2} asks how bootstrap placement interacts with depth and compact state. \textbf{RQ3} asks what CUDA execution and its constituent optimizations contribute to measured training cost. A secondary continuity experiment revisits the original Iris setting; it is an interactive demonstration, not evidence of non-interactive training performance.

Our contributions are: (i) a rectangular, batched diagonal layout with a rotation-only transpose and explicit reduction costs; (ii) a compact-state non-interactive update whose operation counts can be checked against structural predictions; and (iii) an implementation and evaluation protocol that separates cryptographic assumptions, counter-derived comparisons, measured timing, and missing runs. Hoisting, baby-step giant-step (BSGS), sparse-secret encapsulation (SSE), and bootstrapping are established techniques; the contribution lies in their training-specific composition and accounting, not in inventing these primitives.

\paragraph{Status of the evidence.}
This is an interim full draft. The generated tables retain pending rows for the registered full-setting encrypted experiment and GPU training ablation. The completed encrypted subset runs use a different input dimension and cannot replace those rows. The released baseline and our preset also have different secret distributions and security estimates. We consequently do not claim that the registered equal-security, accuracy-matched research goal has been met. Section~\ref{sec:limits} identifies what would change that conclusion.

\section{Preliminaries and Threat Model}
\label{sec:prelim}
\paragraph{Arithmetic.}
CKKS~\cite{ckks2017} encodes approximate vectors in
$R_Q=\mathbb{Z}_Q[X]/(X^N+1)$, with up to $N/2$ complex slots. We use real slot values. A ciphertext at level $\ell$ uses $q_0,\ldots,q_\ell$ and a scale $\Delta$. Multiplication followed by rescaling consumes a modulus prime. Rotations permute slots through a Galois automorphism and a key switch. Hybrid RNS key switching uses special primes $P$ and a digit decomposition; its cost depends on the active limbs, digit sizes, and memory traffic.

We count rotations (including conjugations), relinearizations, and explicit SSE switches:
\begin{equation}
  \KS = n_{\mathrm{rot}}+n_{\mathrm{relin}}+n_{\mathrm{SSE}}.
  \label{eq:ks}
\end{equation}
A ciphertext tensor product is not automatically a key switch: lazy relinearization may combine several products before one switch. Bootstrap events are recorded separately and their constituent operations are already in Eq.~\eqref{eq:ks}; adding a bootstrap-equivalent charge again would double count. Reported counts cover the whole network and are averaged over processed examples, not over layers. A count is a workload metric, not a machine-independent latency estimate.

\paragraph{Refreshes.}
An interactive refresh decrypts and re-encrypts at the key owner. A bootstrap instead evaluates a refresh circuit using public evaluation keys. Our non-interactive path uses SSE~\cite{sse2022}, centered modulus raising, coefficient-to-slot and slot-to-coefficient transforms, and an approximate modular-reduction polynomial. Public transforms use collapsed stages and BSGS; exact scale tracking makes them composable with training.

\paragraph{Protected objects and public leakage.}
We consider a semi-honest evaluation server and a single trusted key owner, under the usual RLWE and evaluation-key assumptions of the chosen CKKS construction. No new security reduction is claimed. Training examples, labels, model weights, intermediate activations, gradients, and optimizer state are encrypted. The architecture, dimensions, batch schedule, hyperparameters, parameters, evaluation keys, and algorithmic control flow are public. The key owner initializes and encrypts the data and initial model, and decrypts the resulting model. Known initialization distributions are not themselves secret. This model does not hide dimensions, resource use, access schedules, or timing, and it does not provide malicious-server verification, differential privacy, multi-party key custody, or protection against inference from a legitimately decrypted model.

\paragraph{Diagnostics are not a server capability.}
The benchmark process can also instantiate a secret-key holder, decrypt checkpoints and predictions, compare against a plaintext twin, and report accuracy or activation ranges. Those are client-side diagnostics for a trusted experiment harness; they must not be exposed as server logs in a deployment. The generic non-interactive trainer has no secret key and makes public depth-based refresh decisions. Disabling diagnostic decryption changes neither the encrypted arithmetic nor the public schedule. The interactive path explicitly involves the client during training and is labeled separately. The in-library ReBoot-layout microbenchmark uses plaintext labels and client-side optimizer updates: it supports arithmetic/count comparisons, not the full confidentiality claim or a faithful non-interactive baseline.

\section{Rotation-Only Transposition for Mini-Batch Training}
\label{sec:packing}
\subsection{A common layout for activations and weights}
Let $B$ be a power-of-two lane count and $d$ a power-of-two feature capacity such that $dB$ divides the slot count. Slot $jB+b$ stores feature $j$ of sample $b$. Blocks repeat when $dB$ is smaller than the slot count. A layer has positive integer input and output dimensions $d_i,d_o$. Define the padded dimensions $n_i=2^{\ceil{\log_2 d_i}}$ and $n_o=2^{\ceil{\log_2 d_o}}$, and require $n_i,n_o\le d$. A vector of padded dimension $n\in\{n_i,n_o\}$ is canonical when its feature sequence is $n$-periodic in every lane; unused logical features are zero. This periodicity eliminates a separate replication step at tall layers.

Zero-pad the original $d_o\times d_i$ weight matrix to $n_o\times n_i$. Let $m=\min(n_o,n_i)$ and $L=\max(n_o,n_i)$. Both divide $d$, so the $L$-periodic diagonals, rotations modulo $LB$, and integer logarithmic folds below are well-defined. Store the generalized diagonals
\begin{equation}
 D_c[j]=W[j\bmod n_o,(j+c)\bmod n_i],\quad 0\le c<m.
 \label{eq:diagonal}
\end{equation}
Weights are shared across all lanes. With BSGS split $m=gu$ and $c=ga+b$, the stored diagonal is
$P_c=\rot(D_c,-gaB)$. The encrypted product is
\begin{equation}
 y=\sum_{a<u}\rot\!\left(\sum_{b<g}P_{ga+b}\odot\rot(x,bB),gaB\right).
 \label{eq:bsgs}
\end{equation}
Here a positive rotation reads slot $s+k$ into output slot $s$; $\odot$ denotes a slot-wise product. For a wide layer, $\log_2(n_i/n_o)$ rotations and additions fold partial outputs together. For a tall layer, the canonical input already has the required replication. Neither case adds a mask to the fold.

\subsection{Transpose without products or masks}
Rectangular diagonal transposition without ciphertext multiplications is inherited from Mihara et al.~\cite{mihara2020}. Their general method uses rotations and additions; their implementation also reports OpenMP batch processing. Here we compose that idea with the periodic $m$-diagonal representation, feature-major SIMD lanes, and BSGS pre-rotations.

The transpose diagonal $E_i[j]=W[(j+i)\bmod n_o,j\bmod n_i]$ equals a rotation of one $D_c$. Choose $c=(-i)\bmod m$, and set $t=-c$ for a wide layer or $t=i$ for a tall layer. In stored BSGS form,
\begin{equation}
 Q_i=\rot(P_c,\,[t-g\floor{i/g}+g\floor{c/g}]B).
 \label{eq:transpose}
\end{equation}
The rotation amount is taken modulo $LB$. Thus transposition uses neither a ciphertext product nor a mask and consumes no multiplicative level. It is \emph{not} free: each nonzero rotation requires a key switch. Dedicated keys minimize rotation count; composing existing baby/giant keys saves key memory at the expense of extra switches. Our training runs use key reuse, and the extra work is included in the counters.

\subsection{Gradients, reductions, and a symbolic cost model}
For layer gradient $\delta$ and input $a$, a pre-rotated weight-gradient diagonal is
\begin{equation}
 \nabla P_{ga+b}
 =\operatorname{batchsum}\bigl(\rot(\delta,-gaB)\odot\rot(a,bB)\bigr).
 \label{eq:gradient}
\end{equation}
Batch summation first reduces lanes by powers-of-two rotations. A lane-zero mask isolates the reduced value; inverse-direction rotations replicate it when replicated gradients are required. The mask has nonintegral polynomial coefficients in the CKKS embedding, so encoding it at unit scale is not an exact, zero-depth shortcut. It consumes a level. Constants for the learning rate and batch normalization are folded into a necessary plaintext multiplication at its exact target scale.

\begin{table}[t]
\centering\small
\caption{Symbolic layer costs for the implemented layout, before any wide-layer fold. $g,u$ are the BSGS split, $m=gu$. Gradient costs show the replicated interactive form; compact non-interactive state changes the final replication.}
\label{tab:symbolic}
\begin{tabular}{lcc}
\toprule
Operation & Rotations & Relins \\
\midrule
BSGS encrypted matvec & $g+u-2$ & $u$ \\
Dedicated transpose & $m-1$ & $0$ \\
Replicated weight gradient & $g+u-2+2m\log_2 B$ & $m$ \\
\bottomrule
\end{tabular}
\end{table}
Hoisting reuses decomposition for rotations of the same input~\cite{halevi2018}. It changes wall-clock work but not Eq.~\eqref{eq:ks}. Lazy relinearization applies within a matvec giant-step sum. It cannot remove the separate relinearizations needed before gradient-diagonal rotations. These distinctions matter when interpreting ablation results.

\paragraph{Extensions relative to the interactive design.}
The first two extensions studied here are the periodic $\min(n_o,n_i)$-diagonal representation with feature-major SIMD mini-batch lanes, and its composition with hoisting and BSGS. General rectangular transposition and batch processing themselves are inherited capabilities, not novelty claims. The remaining four extensions---bootstrap placement, polynomial activations, momentum, and loss conventions---are specified in Section~\ref{sec:noninteractive}. None removes the cost of maintaining canonical server-side state; the initial interactive cost model omitted some replication work and was therefore too optimistic.

\section{Non-Interactive Training}
\label{sec:noninteractive}
\subsection{Compact state and the complete update}
Refreshing every replicated diagonal separately wastes bootstrap slots. We instead reduce each gradient to one lane per feature and combine masking, scaling, and lane shifting in a plaintext product. Selected lanes from several diagonals form compact weight and momentum containers. A public compact-layout plan fixes their offsets and the required rotations; unpacking restores the periodic diagonals for the next forward pass. Bias slots can share unused compact positions. Masked packing and unpacking both have a depth cost and are included in the closed-form predictor.

The server performs encrypted forward propagation, forms the loss gradient from encrypted labels, propagates through the transposed weights, and accumulates gradients over micro-batches before one optimizer update. Increasing the effective batch beyond the physical lane count therefore amortizes optimizer refreshes without pretending the examples fit in one ciphertext. It also changes optimization behavior and cannot be treated as a cost-only substitution.

The original dense-secret bootstrap constants did not satisfy the desired precision and depth budget. The implemented preset instead uses a dense ternary main key and a temporary sparse key during modulus raising. The SSE pair is returned to the dense main key before the remaining arithmetic. Collapsed transforms and a level-optimal Chebyshev evaluator fit the selected modulus bands. This is an engineering parameter choice, not a new bootstrap construction or an inherited failure-probability theorem. Appendix~\ref{app:params} specifies the chain and the exposed key instances.

\subsection{Three refresh schedules (extension 3)}
\texttt{AfterUpdate} (AU) bootstraps every compact weight and velocity container after each update. \texttt{WhenNeeded} (WN), an additional schedule beyond the original plan, does so only if the next update would exceed the remaining depth. The decision uses the next batch's fullness: padded lanes at a linear output need a validity mask, so a partial batch can consume an extra level. An unknown next shape is conservatively treated as partial.

\texttt{AfterActivation} (AA) refreshes each hidden activation and retains the WN rule for compact weights. For the shallow eMLP-1 workload, activation refresh does not reduce the frequency of weight refresh beyond WN; it adds a bootstrap per hidden activation and micro-batch. Its motivation is enabling deeper circuits that otherwise exceed the budget, not an automatic speedup. Activation magnitudes must fit the supported bootstrap range. The observed AA subset run exceeded the conservative range in some slots and had a larger activation error; it was not clipped, and does not certify arbitrary activation distributions.

For a sequence of identical updates consuming $D$ levels with $L_a$ usable levels after unpacking, $\floor{L_a/D}$ is only an initial scheduling bound. Unpack costs between updates, partial batches, and the actual state level must also enter the scheduler. We check the structural predictor against every measured update rather than replacing this stateful schedule with a fixed average.

\subsection{Activation, optimizer, and loss (extensions 4--6)}
\paragraph{Polynomial activation.}
We retain the earlier squared activation for continuity and implement $z^2+z$ for the ReBoot comparison. A polynomial quartic option is available for depth experiments. Its additional depth and available tests do not establish a measured accuracy improvement; that planned claim remains unsupported.

\paragraph{Momentum.}
Both SGD and Nesterov-style momentum are supported. Writing $h=\eta(1+\mu)G/n$ and scaled velocity $U=\eta(1+\mu)V$ gives
\begin{align}
 U'&=\mu U+h,\nonumber\\
 W'&=W-h-\frac{\mu^2}{1+\mu}U.
 \label{eq:nag}
\end{align}
Constants are folded into exact-scale products already needed by the gradient path. A changing learning rate requires rescaling the stored velocity representation, and its lane shifts are charged. Biases use SGD; the faithful ReBoot configuration freezes their encrypted zero values while retaining the implementation's bias arithmetic. Weight decay is not implemented; the captured faithful setting uses zero decay.

\paragraph{Loss.}
A squared-error loss yields an affine output gradient; MSE normalization and ReBoot's sum reduction are matched through the learning-rate convention. The faithful configuration uses the released initialization, sample order, output padding, and schedule, not just the paper's nominal hyperparameter names. Softmax/cross-entropy approximation is optional in the plan and is not implemented or evaluated here. Block MLPs follow local classifier losses and stop gradients across blocks; they must not be described as an ordinary deeper back-propagation network.

\section{GPU Implementation}
\label{sec:gpu}
\sys{} separates RNS arithmetic, CKKS evaluation, neural-network operations, CUDA execution, and benchmark orchestration. Generic \texttt{HeOps} and \texttt{BootstrapOps} interfaces let the same trainer use CPU or device ciphertexts. The evaluator checks parameter identity, levels, scales, key availability, and device ownership before work. Counters commit after success; the CUDA bootstrap also waits for final device synchronization before merging its internal counters.

The CUDA backend uses device NTTs, basis conversion, pointwise arithmetic, hybrid key switching, and BSGS matvecs. Small public launch data use pinned transfer buffers. Reusing a hoisted decomposition avoids repeating the expensive preparation for related rotations. Device/CPU bit parity is meaningful only for cases actually executed; memory-skipped test functions can still report process success.

\paragraph{Streams and memory.}
Independent products may be issued to several streams, but a training step's dependent matvecs are executed sequentially. The measured pre-fix multi-stream benchmark was slower than sequential execution (Table~\ref{tab:gpu_vs_cpu_ops}). Calling-thread submission and pool handling were subsequently revised; their performance remeasurement is pending. An effective-bandwidth estimate for addition does not prove matvec occupancy or full-device saturation. That explanation remains a hypothesis until profiled.

The shared device has limited available memory. Training can upload keys capped to an explicitly selected level; this preserves the computation at that level, but does not measure the top-level workload. Allocation preflight uses estimated resident/key/workspace sizes. Pool trims occur after large setup stages, at epoch boundaries, and between ablation configurations. Estimated need and pool checkpoints are not a measured memory peak, and a successful preflight does not guarantee later allocations under contention.

\paragraph{Bootstrap boundary.}
The CUDA bootstrap implementation keeps ciphertext arithmetic on the device through SSE, modulus raising, transforms, and EvalMod. A narrow shared polynomial executor preserves the CPU operation order and rounding. Public integer constants are still materialized on the host and uploaded; future device timing must include that overhead. Local CPU replay and CUDA compilation support code readiness, not physical-GPU precision, latency, or memory feasibility. Those hardware measurements are pending.

The training benchmark still uses a CPU bootstrap adapter. Level operations run on the GPU; non-interactive refresh downloads to the CPU and uploads the result. A compiled device bootstrap API has not changed that benchmark. Its timing must therefore separate device work, transfers, and CPU refresh. Likewise interactive GPU training uses a client refresh. Neither path justifies an end-to-end all-GPU training claim.

\section{Evaluation}
\label{sec:evaluation}
\subsection{Protocol and evidence classes}
We distinguish a released baseline execution, an in-library layout microbenchmark, and a cost model. The first supplies observed wall-clock and bootstrap events; the second isolates arithmetic layout counts under our evaluator; the third supplies uninstrumented baseline quantities. A counter-derived normalization is not a separate execution. Historical generated tables retain the working name \emph{hekit} for \sys{} and its crate names.

The headline uses MNIST reduced to $196$ inputs and the released full training setting. FULL-1 has completed at $B=64$ with one initial-weight seed; FULL-2 at the released baseline batch remains pending. FULL-1 is partial full-setting evidence, not a substitute for the FULL-2 accuracy gate. The completed faithful encrypted subset uses $784$ inputs, the baseline's sample indices and initial weights, and its training conventions. Output width is padded to $16$, including the unused classes in the released loss; replacing this with a width-$10$ loss is not assumed equivalent. Earlier ``legacy'' runs use a different learning rate, bias policy, and initialization. Their counters and timing remain useful, but their accuracy is not comparable to the released baseline. The last subset batch is partial, affecting depth, refresh scheduling, and the effective sum-versus-mean learning-rate conversion.

The paired server measurements use the same CPU and thread cap, with different CPU parallel runtimes. Jobs are not core-pinned, and some rows overlap unrelated shared-host load. Client work and setup scopes also differ: baseline step time includes batch encryption and prediction decryption, whereas our update time excludes those separately reported client operations; both exclude key setup. Thus even a same-host ratio is not a controlled architectural speedup. We do not pool CPU and GPU samples. Encrypted rows have a single seed unless the table states otherwise; within-run iterations and Criterion samples are not independent seeds.

\subsection{RQ1: key switches and wall-clock}
% generated by tools/paper_tables.py from results/; do not edit
\begin{table*}[t]
\centering
\scriptsize
\caption{RQ1: key switches per sample, four ways (ratio hekit / ReBoot at B = 48)}
\label{tab:rq1_ks}
\resizebox{\linewidth}{!}{%
\begin{tabular}{lrrrrrrrrrrrrr}
\toprule
Run & Input & Placement & B & Config & Samples & hekit KS/sample & Boots & ReBoot non-boot & (a) & (b) & (c) & (d) & Ways $\leq$ 0.70 \\
\midrule
\multicolumn{14}{l}{\textit{196 inputs, ReBoot full setting, 4800 samples x 10 epochs, seed 0 (headline)}} \\
FULL-1 & 196 & WN & 64 & faithful & 48000 & 18.272 & 1500 & 37 (model) & 0.232 & 0.387 & 0.395 & 0.400 & abcd \\
FULL-2 & 196 & WN & 48 in 64 & faithful & 48000 & 28.990 & 4000 & 37 (model) & 0.368 & 0.781 & 0.627 & 0.534 & acd \\
\midrule
\multicolumn{14}{l}{\textit{784 inputs, ReBoot encrypted-script subset, 1 epoch (supplementary)}} \\
F1 & 784 & WN & 32 & faithful & 1008 & 33.488 & 68 & 35.0625 (Lazy, measured) & 0.435 & 0.779 & 0.756 & 0.742 & a \\
F2 & 784 & WN & 64 & faithful & 1008 & 17.885 & 36 & 35.0625 (Lazy, measured) & 0.233 & 0.414 & 0.404 & 0.397 & abcd \\
R-41-3 & 784 & WN & 32 & legacy (R-M3-12) & 1024 & 32.531 & 64 & 35.0625 (Lazy, measured) & 0.423 & 0.741 & 0.734 & 0.730 & a \\
R-41-4 & 784 & WN & 64 & legacy (R-M3-12) & 1024 & 17.172 & 32 & 35.0625 (Lazy, measured) & 0.223 & 0.382 & 0.388 & 0.391 & abcd \\
R-41-1 & 784 & AU & 32 & legacy (R-M3-12) & 1024 & 39.469 & 128 & 35.0625 (Lazy, measured) & 0.513 & 1.149 & 0.891 & 0.730 & a \\
R-41-2 & 784 & AU & 64 & legacy (R-M3-12) & 1024 & 20.641 & 64 & 35.0625 (Lazy, measured) & 0.268 & 0.586 & 0.466 & 0.391 & abcd \\
R-41-5 & 784 & AA & 32 & legacy (R-M3-12) & 1024 & 36.000 & 96 & 35.0625 (Lazy, measured) & 0.468 & 0.945 & 0.812 & 0.730 & a \\
\midrule
\multicolumn{14}{l}{\textit{Interactive hekit, client refresh, no bootstraps (separate row, not way (d))}} \\
interactive & 784 & --- & 32 & legacy (M2) & 1024 & 19.781 & 0 & 35.0625 (Lazy, measured) & --- & --- & --- & --- & --- \\
\bottomrule
\end{tabular}}
\par\smallskip\begin{minipage}{\linewidth}\scriptsize
Independent interactive reference: ratio 0.564 against ReBoot Lazy; client refresh, not NI-derived way (d) and not a four-way pass result.\par
KS = rotations + relinearizations + SSE key switches (OpCounter). hekit bootstrap = 111 KS (measured, setup.boot\_ks). ReBoot bootstrap = 502 KS and 196-input non-boot = 37 KS/sample are M0 cost-model values (OpenFHE does not count key switches); ReBoot bootstraps 4 ciphertexts per update at B = 48 (measured on the official subset run: 84 bootstraps / 21 steps).\par
(a) own bootstrap costs; (b) both 502; (c) both 111; (d) bootstraps excluded on both sides. (b)-(d) are derived from the measured counters, not separate executions. Criterion: ratio $\leq$ 0.70 (research plan \S{}4.3.2).\par
784 rows are the supplementary subset rows (R-M3-13): ReBoot non-boot = Lazy 35.06, the in-library measurement of ReBoot's layout (Task 33, baseline c). Legacy rows use the pre-R-M3-12 config; their counts equal the faithful config but their accuracy is not comparable.\par
\end{minipage}
\end{table*}

% generated by tools/paper_tables.py from results/; do not edit
\begin{table*}[t]
\centering
\scriptsize
\caption{Wall-clock on the same CPU (quadruped-gpu, i9-13900KF, 12 threads)}
\label{tab:wallclock_same_cpu}
\resizebox{\linewidth}{!}{%
\begin{tabular}{lrrrrrrrrr}
\toprule
Run & Input & B & Samples & Train s & s/sample & vs ReBoot & Boot s/sample & Non-boot s/sample & Threads \\
\midrule
\multicolumn{10}{l}{\textit{196 inputs, full setting (headline)}} \\
FULL-1 & 196 & 64 & 48000 & pending (R-41-FULL-1) &  & ReBoot 196 full: not measured &  &  &  \\
FULL-2 & 196 & 48 in 64 & 48000 & pending (R-41-FULL-2) &  & ReBoot 196 full: not measured &  &  &  \\
\midrule
\multicolumn{10}{l}{\textit{784 inputs, ReBoot encrypted-script subset}} \\
ReBoot official (OpenFHE) & 784 & 48 & 1008 & 3599.9 & 3.571 & 1 & 2.585 & 0.986 & OpenMP 12 \\
F1 & 784 & 32 & 1008 & 3016.7 & 2.993 & 0.838 & 0.989 & 2.003 & rayon 12 \\
F2 & 784 & 64 & 1008 & 1607.9 & 1.595 & 0.447 & 0.522 & 1.073 & rayon 12 \\
R-41-3 & 784 & 32 & 1024 & 2773.6 & 2.709 & 0.758 & 0.871 & 1.838 & rayon 12 \\
R-41-4 & 784 & 64 & 1024 & 1434.8 & 1.401 & 0.392 & 0.435 & 0.967 & rayon 12 \\
R-41-1 & 784 & 32 & 1024 & 4339.1 & 4.237 & 1.187 & 1.736 & 2.501 & rayon 12 \\
R-41-2 & 784 & 64 & 1024 & 2228.9 & 2.177 & 0.609 & 0.867 & 1.310 & rayon 12 \\
R-41-5 & 784 & 32 & 1024 & 3606.6 & 3.522 & 0.986 & 1.316 & 2.206 & rayon 12 \\
\bottomrule
\end{tabular}}
\par\smallskip\begin{minipage}{\linewidth}\scriptsize
ReBoot time = sum of its own step\_time (encrypt batch, forward/backward, decrypt prediction, weight bootstrap; setup excluded). hekit time = time\_s.train\_total (all updates including bootstraps; key generation and bootstrapper setup excluded). Present rows fail closed unless the binary's provenance names the same host and thread count as the ReBoot run.\par
Shared machine: other users' jobs (SLAM, \textasciitilde{}4 cores) ran during some rows; see docs/m3-report.md \S{}3.2. ReBoot's full 196-input setting was not run (estimated 47.6 h), so the headline has no same-CPU wall-clock ratio.\par
\end{minipage}
\end{table*}

Let $K$ be measured total key switches, $H$ bootstrap events, and $S$ examples processed. Let $c_h$ be our measured per-bootstrap cost and $c_r$ the baseline bootstrap model; let $R$ denote the baseline non-bootstrap count per example, $r$ its refreshes per update, and $B_r$ its own batch. The four reported ratios are
\begin{align}
 \rho_a&=\frac{K/S}{R+rc_r/B_r},&
 \rho_b&=\frac{(K-c_hH+c_rH)/S}{R+rc_r/B_r},\nonumber\\
 \rho_c&=\frac{K/S}{R+rc_h/B_r},&
 \rho_d&=\frac{(K-c_hH)/S}{R}.
 \label{eq:ratios}
\end{align}
Convention (a) uses each side's own bootstrap cost. Convention (b) charges both the baseline cost; (c) charges both our cost; (d) excludes bootstrap arithmetic on both sides. The interactive comparison is separate from (d): omitting bootstrap arithmetic does not remove non-interactive compact packing and unpacking. ReBoot remains at its own batch in the denominator; a hypothetical same-batch denominator would answer a different question.

The baseline's bootstrap cost remains a model because the measured OpenFHE execution does not expose key-switch counters. Its observed bootstrap frequency is used, not replaced with our schedule. The subset non-bootstrap reference is the in-library Lazy layout measurement; the full-setting reference is still a model. These boundaries are visible in Table~\ref{tab:rq1_ks}.

% Generated by tools/figs.py; do not edit. See rq1_ratios.sources.json.
\begin{figure}[t]
\centering
\setlength{\unitlength}{1pt}
\begin{picture}(240,143)
\put(28,32){\line(1,0){204}}
\put(28,32){\line(0,1){82}}
\put(5,124){\makebox(225,10){\footnotesize KS ratio to ReBoot (lower is better)}}
\put(25,32.00){\line(1,0){3}}
\put(0,29.00){\makebox(22,7)[r]{\tiny 0}}
\put(25,52.00){\line(1,0){3}}
\put(0,49.00){\makebox(22,7)[r]{\tiny 0.25}}
\put(25,72.00){\line(1,0){3}}
\put(0,69.00){\makebox(22,7)[r]{\tiny 0.5}}
\put(25,92.00){\line(1,0){3}}
\put(0,89.00){\makebox(22,7)[r]{\tiny 0.75}}
\put(25,112.00){\line(1,0){3}}
\put(0,109.00){\makebox(22,7)[r]{\tiny 1}}
\multiput(28,88.00)(5,0){41}{\line(1,0){2}}
\put(176,91.00){\makebox(55,7)[r]{\tiny threshold 0.7}}
\put(42,32){\color[HTML]{315C8C}\rule{14pt}{34.800pt}}
\put(38,68.800){\makebox(22,7){\tiny 0.435}}
\put(42,21){\makebox(14,7){\scriptsize (a)}}
\put(61,32){\color[HTML]{487F9C}\rule{14pt}{62.320pt}}
\put(57,96.320){\makebox(22,7){\tiny 0.779}}
\put(61,21){\makebox(14,7){\scriptsize (b)}}
\put(80,32){\color[HTML]{6C9B8B}\rule{14pt}{60.480pt}}
\put(76,94.480){\makebox(22,7){\tiny 0.756}}
\put(80,21){\makebox(14,7){\scriptsize (c)}}
\put(99,32){\color[HTML]{A2A76B}\rule{14pt}{59.360pt}}
\put(95,93.360){\makebox(22,7){\tiny 0.742}}
\put(99,21){\makebox(14,7){\scriptsize (d)}}
\put(35,6){\makebox(90,10){\footnotesize F1: B = 32}}
\put(145,32){\color[HTML]{315C8C}\rule{14pt}{18.640pt}}
\put(141,52.640){\makebox(22,7){\tiny 0.233}}
\put(145,21){\makebox(14,7){\scriptsize (a)}}
\put(164,32){\color[HTML]{487F9C}\rule{14pt}{33.120pt}}
\put(160,67.120){\makebox(22,7){\tiny 0.414}}
\put(164,21){\makebox(14,7){\scriptsize (b)}}
\put(183,32){\color[HTML]{6C9B8B}\rule{14pt}{32.320pt}}
\put(179,66.320){\makebox(22,7){\tiny 0.404}}
\put(183,21){\makebox(14,7){\scriptsize (c)}}
\put(202,32){\color[HTML]{A2A76B}\rule{14pt}{31.760pt}}
\put(198,65.760){\makebox(22,7){\tiny 0.397}}
\put(202,21){\makebox(14,7){\scriptsize (d)}}
\put(138,6){\makebox(90,10){\footnotesize F2: B = 64}}
\end{picture}
\caption{Supplementary subset RQ1 ratios (one seed per row), generated from the same cells as Table~\ref{tab:rq1_ks}. Bars (a)--(d) denote the four accounting conventions in Section~\ref{sec:evaluation}. These are not equal-security or full-setting results.}
\Description{Grouped bars for the two faithful subset batches under four key-switch accounting conventions. The larger batch is below the threshold in all four; the smaller batch only under own bootstrap costs. One seed per row; not an equal-security comparison.}
\label{fig:rq1-ratios}
\end{figure}

The faithful larger-batch subset row meets the registered count threshold under every reported convention, while the smaller batch meets it only under the own-bootstrap convention. The generated figure makes that dependence visible. The corresponding larger-batch timing is also favorable on the shared CPU (Table~\ref{tab:wallclock_same_cpu}), but its subset accuracy drops relative to the baseline, including in the plaintext twin. Consequently this observation is a tradeoff, not a successful accuracy-matched RQ1 result. FULL-1 also meets the count threshold under all four conventions. Its completed training time has no measured full-setting ReBoot CPU denominator. The pending FULL-2 accuracy gate, missing baseline timing, and security mismatch below prevent a final verdict.

\subsection{RQ2: placement and depth}
The AU, WN, and AA subset rows in Tables~\ref{tab:rq1_ks} and~\ref{tab:wallclock_same_cpu} provide the placement comparison. For matched legacy eMLP-1 settings, WN reduces bootstrap events and total time relative to AU. AA leaves the same lazy weight-refresh frequency and adds activation refreshes, so it does not improve this shallow workload. Its benefit is a feasible depth budget for some deeper polynomial circuits, checked separately as a structural test rather than an observed training-accuracy gain.

The corrected structural predictor equals the measured counters at every completed update. This validates accounting for the executed paths; it does not validate the baseline bootstrap model or a convergence claim. The initial cost model missed replication, key-reuse, bias, and schedule-related terms, and its optimistic interactive estimate is superseded by that explicit accounting. FULL-1 extends exact schedule agreement to its complete registered update sequence and records accumulated weight and bias error against its same-batch plaintext twin. That observed error is larger than in the earlier subset checks; close final accuracy does not imply a uniform small parameter error. FULL-2 and guarantees across seeds or arbitrary trajectories remain unmeasured.

\subsection{Accuracy, local blocks, and security}
% generated by tools/paper_tables.py from results/; do not edit
\begin{table*}[t]
\centering
\scriptsize
\caption{Test accuracy against ReBoot (R-M3-15)}
\label{tab:accuracy_vs_reboot}
\resizebox{\linewidth}{!}{%
\begin{tabular}{lrrrrrrr}
\toprule
Row & Setting & Seeds & Test acc (\%) & Twin (\%) & Reference (\%) & $\Delta$ (pt) & Gate \\
\midrule
\multicolumn{8}{l}{\textit{Full setting: 196 inputs, 4800 samples x 10 epochs, B = 48 (the $\pm$0.5 pt gate)}} \\
ReBoot released code, plaintext (recomputed twin) & full & 10 & 90.65 $\pm$ 0.27 (n = 10) & --- & --- & --- & reference \\
hekit plaintext twin (ReBoot weights, order, lr) & full & 10 & 90.65 $\pm$ 0.27 (n = 10) & --- & 90.65 $\pm$ 0.27 (n = 10) & +0.00 & twin check \\
hekit encrypted, B = 48 in 64 lanes, WN & full & 1 (seed 0) & pending (R-41-FULL-2) & pending (R-41-FULL-2) & 90.41 &  & pending \\
hekit encrypted, B = 64, WN & full & 1 (seed 0) & 90.36 & 90.37 & 90.41 (seed 0) & -0.05 & not the gate (B $\neq$ 48) \\
\midrule
\multicolumn{8}{l}{\textit{eMLP-2 / eMLP-3 block MLPs, full setting, seed 0}} \\
hekit plaintext twin emlp2 (blocks 64) & full & 1 (seed 0) & 91.56 & --- & 91.56 (ReBoot twin) & +0.00 & twin check \\
hekit plaintext twin emlp3 (blocks 128,64) & full & 1 (seed 0) & 91.75 & --- & 91.75 (ReBoot twin) & +0.00 & twin check \\
\midrule
\multicolumn{8}{l}{\textit{784 inputs, ReBoot encrypted-script subset, 1008 samples, 1 epoch (not a gate, R-M3-15)}} \\
ReBoot official encrypted (OpenFHE) & subset, B = 48 & 1 (seed 0) & 75.06 & --- & --- & --- & reference \\
hekit encrypted WN, B = 32 & subset & 1 (seed 0) & 75.77 & 75.77 & 75.06 & +0.71 & not a gate \\
hekit encrypted WN, B = 64 & subset & 1 (seed 0) & 73.58 & 73.59 & 75.06 & -1.48 & not a gate \\
hekit encrypted eMLP-2, B = 48 (196) & subset 1008 & 1 & pending (R-52-1) & pending (R-52-1) &  &  & twin check \\
hekit encrypted eMLP-3, B = 48 (196) & subset 1008 & 1 & pending (R-52-2) & pending (R-52-2) &  &  & twin check \\
\midrule
\multicolumn{8}{l}{\textit{Reference only: the paper value and the released code at the paper setting (R-REB-3)}} \\
ReBoot paper, Table 4 (plaintext) & 60k samples & 10 (paper) & 94.61 $\pm$ 0.27 & --- & --- & --- & not reproducible from the released defaults \\
ReBoot released code, paper setting, constant lr, epoch 50 & 60k samples, 50 epochs & 10 & 95.15 $\pm$ 0.19 (n = 10) & --- & 94.61 & +0.54 & reference \\
ReBoot released code, paper setting, constant lr, best epoch (test-selected) & 60k samples, 50 epochs & 10 & 95.35 $\pm$ 0.10 (n = 10) & --- & 94.61 & +0.74 & reference \\
ReBoot released code, paper setting, ReBoot's scheduler, epoch 50 & 60k samples, 50 epochs & 10 & 95.25 $\pm$ 0.08 (n = 10) & --- & 94.61 & +0.64 & reference \\
ReBoot released code, paper setting, ReBoot's scheduler, best epoch (test-selected) & 60k samples, 50 epochs & 10 & 95.46 $\pm$ 0.09 (n = 10) & --- & 94.61 & +0.85 & reference \\
\bottomrule
\end{tabular}}
\par\smallskip\begin{minipage}{\linewidth}\scriptsize
FULL-1 final parameter error: 0.00682657 (maximum absolute difference across weights and biases against its same-batch plaintext twin at training end). This is a trajectory diagnostic, not a per-bootstrap precision or accuracy gate.\par
R-M3-15: the $\pm$0.5 pt gate compares hekit encrypted B = 48 (FULL-2) with ReBoot's released code at its full setting, seed 0. Subset rows are 1 epoch and are not a gate. The paper value and R-REB-3 are reported with this note only.\par
\end{minipage}
\end{table*}

% generated by tools/paper_tables.py from results/; do not edit
\begin{table*}[t]
\centering
\scriptsize
\caption{Plaintext dataset reproduction against ReBoot Table 4 (Task 51)}
\label{tab:plain_datasets}
\resizebox{\linewidth}{!}{%
\begin{tabular}{lrrrrrrrr}
\toprule
Dataset & Config & Epochs & Seeds & Fixed final epoch (\%) & Best epoch, test-selected (\%) & ReBoot Table 4 (\%) & $\Delta$ fixed (pt) & Within 1 pt \\
\midrule
MNIST 14x14 & paper setting & 50 & 10 & 95.00 $\pm$ 0.18 (n = 10) & 95.33 $\pm$ 0.15 (n = 10) & 94.61 $\pm$ 0.27 & +0.39 & yes \\
Fashion-MNIST 14x14 & paper setting & 50 & 10 & 85.22 $\pm$ 0.17 (n = 10) & 85.53 $\pm$ 0.12 (n = 10) & 85.12 $\pm$ 0.20 & +0.10 & yes \\
Breast Cancer & fixed split, min-max (primary) & 100 & 10 & 96.40 $\pm$ 1.27 (n = 10) & 99.30 $\pm$ 0.37 (n = 10) & 99.03 $\pm$ 1.78 & -2.63 & no \\
Breast Cancer & per-run split (sensitivity) & 100 & 10 & 95.26 $\pm$ 1.61 (n = 10) & 97.89 $\pm$ 0.94 (n = 10) & 99.03 $\pm$ 1.78 & -3.77 & no \\
Breast Cancer & z-score (diverges) & 100 & 10 & diverged (10/10 seeds) & --- & 99.03 $\pm$ 1.78 & --- & no \\
\bottomrule
\end{tabular}}
\par\smallskip\begin{minipage}{\linewidth}\scriptsize
Fixed final epoch is the primary figure. The best epoch is selected on the test set and is optimistic; it is shown only because ReBoot does not state its epoch count. Seeds: the plain binary's own seeds (not ReBoot's initial weights).\par
Fashion-MNIST and Breast Cancer are plaintext twins only (no encrypted runs).\par
\end{minipage}
\end{table*}

% generated by tools/paper_tables.py from results/; do not edit
\begin{table*}[t]
\centering
\scriptsize
\caption{Security estimates (lattice-estimator; log2 rop, minimum over attacks)}
\label{tab:security}
\resizebox{\linewidth}{!}{%
\begin{tabular}{lrrrrrr}
\toprule
Parameters & N & log PQ & Secret & BDGL16 & MATZOV & >= 128 (both) \\
\midrule
hekit (a): n16\_boot main key & 2\textasciicircum{}16 & 1732 (1731.998) & dense ternary (P($\pm$1) = 1/4 each, h $\approx$ N/2) & 131.5 & 129.3 & yes \\
hekit (b): SSE key s -> s' (h' = 32) & 2\textasciicircum{}16 & 294 (q\_0$\cdot$P) & sparse ternary, h' = 32 & 133.8 & 134.2 & yes \\
ReBoot main key (R-REB-4: OpenFHE's own chain) & 2\textasciicircum{}16 & 1730 (1729.585) & sparse ternary, h = 192 (97 / 95) & 121.8 & 121.5 & no \\
\bottomrule
\end{tabular}}
\par\smallskip\begin{minipage}{\linewidth}\scriptsize
ReBoot row: the MITM tier (operation counts; about 2\textasciicircum{}104-2\textasciicircum{}110 memory). Without MITM attacks ReBoot is at 127.8 / 126.2 (BDGL16 / MATZOV), also below 128.\par
Source: the summary table of docs/security.md (the estimator outputs are recorded in that doc; no JSON exists). Same lattice-estimator models for every row; m = N; Core-SVP is in the doc for information only. Equal N does not mean equal security: no equal-security claim.\par
\end{minipage}
\end{table*}

Table~\ref{tab:accuracy_vs_reboot} separates the released full-setting reference, subset results, and published paper accuracy. The released defaults and paper-setting reconstruction are different protocols. We do not choose whichever reference makes the result look best. The registered tolerance is tested by FULL-2 against the captured released full-setting seed-zero reference; that encrypted row has not completed. FULL-1 reports its larger-batch accuracy and same-batch twin separately, without substituting either for this gate. Matching a plaintext twin is necessary to diagnose arithmetic fidelity, but is insufficient to establish baseline-equivalent generalization.

Plaintext dataset experiments (Table~\ref{tab:plain_datasets}) provide broader optimization checks. MNIST and Fashion-MNIST reproduce the stated comparison tolerance at the fixed final epoch. Breast Cancer does not: the fixed-epoch min-max result misses its target, and selecting the best test epoch is optimistic leakage from model selection. The z-score setting diverges for the polynomial activation. Fashion-MNIST and Breast Cancer results are plaintext twins, not encrypted training measurements. The local-block plaintext twins reproduce the released eMLP-2/eMLP-3 values for the recorded seed; their encrypted subset rows are pending. The block experiments do not extend the current headline claim to deeper networks.

Table~\ref{tab:security} applies common estimator models to the actual exposed key moduli and secret distributions. The evaluated main and SSE instances of our preset exceed the chosen operation-count target under those models; the released ReBoot preset does not. This is an estimate over the documented attack search, not a security proof. The smallest sparse-secret values use attacks with extreme memory requirements, and even non-MITM minima have substantial or unreported memory demands. These are not practical attack demonstrations. Equal $N$ and a requested library security label do not establish equal security. Results must remain conditional until either a security-different comparison is explicitly accepted or an independently estimated and remeasured alternative baseline is available.

\subsection{RQ3: device operations and ablation}
% generated by tools/paper_tables.py from results/; do not edit
\begin{table*}[t]
\centering
\scriptsize
\caption{GPU vs CPU vs HEonGPU, n16\_train (R-47-2, criterion medians, ms)}
\label{tab:gpu_vs_cpu_ops}
\resizebox{\linewidth}{!}{%
\begin{tabular}{lrrrrrrr}
\toprule
Operation & Level & hekit GPU & CPU 1 thr & CPU 12 thr & GPU speed-up (1 / 12 thr) & HEonGPU (RTX 4090) & GPU / HEonGPU \\
\midrule
add & 31 & 0.111 & 15.674 & 14.611 & 141x / 132x & none published & --- \\
relinearize & 31 & 4.671 & 639.639 & 240.070 & 137x / 51x & 2.582 (l = 36, I = 4 primes per digit) & 1.81x (nominal) \\
mul + relin + rescale & 31 & 5.823 & 725.209 & 337.731 & 125x / 58x & none published & --- \\
rotate(1) & 31 & 3.990 & 622.046 & 223.880 & 156x / 56x & none published & --- \\
hoisted rotate x16 & 31 & skipped: needs 5580 MiB (2377 MiB free) & 3765.068 & 689.471 & --- & none published & --- \\
hoisted rotate x16 (capped keys) & 9 & 4.984 & 1005.331 & 172.870 & 202x / 35x & none published & --- \\
matvec\_enc d = 64 & 31 & skipped: needs 7276 MiB (1929 MiB free) & 12938.440 & 2325.173 & --- & none published & --- \\
matvec\_enc d = 64 (capped keys) & 9 & 16.120 & 2523.223 & 547.762 & 157x / 34x & none published & --- \\
matvec\_enc d = 64 (capped keys) & 5 & 9.509 & 1391.064 & 314.649 & 146x / 33x & none published & --- \\
matvec\_enc x4, 1 stream (sequential) & 5 & 34.994 & --- & --- & --- & none published & --- \\
matvec\_enc x4, 4 streams (untuned R-47-2) & 5 & 41.673 & --- & --- & --- & none published & --- \\
matvec\_enc x4, re-measured after the fix round & 5 & pending (R-47-3) & --- & --- & --- & --- & --- \\
\bottomrule
\end{tabular}}
\par\smallskip\begin{minipage}{\linewidth}\scriptsize
HEonGPU relinearize: Ozcan and Savas, GPU Implementations of Three Different Key-Switching Methods for Homomorphic Encryption Schemes, IACR ePrint 2025/124, Table 3: CKKS relinearization on a GeForce RTX 4090, log N = 16, 43-bit precision, Method II (hybrid), (I, l\textasciitilde{}, l) = (4, 40, 36), level 35. The previously cited Nov 2024 manuscript version date is unverified; ePrint landing metadata alone does not establish that manuscript date. Caveat: not the same parameters as n16\_train (36 vs 32 Q primes, 4 special primes; I = 4 primes per digit corresponds to alpha = 4, not four digits; n16\_train uses dnum 8); the ratio is nominal. Scaling key switching linearly in l gives about 2.30 ms at l = 32 and a ratio of about 2.0x (docs/m4-report.md section 1.1). ePrint 2024/1543 could not be read (HTTP 403), so other HEonGPU operation values are unconfirmed.\par
Capped keys: the input is dropped to level 9 or 5 and the keys are cut to that level (faithful for that level, not the top-level cost); the CPU uses full keys at the same level. The shared RTX 4090 had 2-3 GiB free, so full-key GPU rows that need more were skipped.\par
\end{minipage}
\end{table*}

% generated by tools/paper_tables.py from results/; do not edit
\begin{table*}[t]
\centering
\scriptsize
\caption{RQ3 ablation: one interactive eMLP-1 iteration under each EvalConfig}
\label{tab:rq3_ablation}
\resizebox{\linewidth}{!}{%
\begin{tabular}{lrrrrrrrrr}
\toprule
Source & Backend / level & hoist bsgs lazy & Rotations & Relins & KS/iter & Closed form & Step median (s) & Step / default & Max |$\Delta$W| vs default \\
\midrule
CPU, n16\_train, level 31 &  &  &  &  &  &  & pending (R-48-4) &  &  \\
GPU, n16\_train, level 7 &  &  &  &  &  &  & pending (R-48-2) &  &  \\
CPU, n16\_train, level 7 &  &  &  &  &  &  & pending (R-48-3) &  &  \\
CPU smoke (insecure N = 2\textasciicircum{}12: counts only) & cpu, level max & 1 1 1 & 171 & 62 & 233 & match & n/a (smoke) & n/a (smoke) & 0 \\
CPU smoke (insecure N = 2\textasciicircum{}12: counts only) & cpu, level max & 1 1 0 & 171 & 114 & 285 & match & n/a (smoke) & n/a (smoke) & 7.83e-09 \\
CPU smoke (insecure N = 2\textasciicircum{}12: counts only) & cpu, level max & 1 0 1 & 210 & 101 & 311 & match & n/a (smoke) & n/a (smoke) & 6.90e-09 \\
CPU smoke (insecure N = 2\textasciicircum{}12: counts only) & cpu, level max & 1 0 0 & 210 & 114 & 324 & match & n/a (smoke) & n/a (smoke) & 7.39e-09 \\
CPU smoke (insecure N = 2\textasciicircum{}12: counts only) & cpu, level max & 0 1 1 & 171 & 62 & 233 & match & n/a (smoke) & n/a (smoke) & 0 \\
CPU smoke (insecure N = 2\textasciicircum{}12: counts only) & cpu, level max & 0 1 0 & 171 & 114 & 285 & match & n/a (smoke) & n/a (smoke) & 7.83e-09 \\
CPU smoke (insecure N = 2\textasciicircum{}12: counts only) & cpu, level max & 0 0 1 & 210 & 101 & 311 & match & n/a (smoke) & n/a (smoke) & 6.90e-09 \\
CPU smoke (insecure N = 2\textasciicircum{}12: counts only) & cpu, level max & 0 0 0 & 210 & 114 & 324 & match & n/a (smoke) & n/a (smoke) & 7.39e-09 \\
\bottomrule
\end{tabular}}
\par\smallskip\begin{minipage}{\linewidth}\scriptsize
1. Definition of bsgs = false. The giant rotation is not shared; the weights keep the same pre-rotated (BSGS) layout and the same Galois keys. Every product with a nonzero giant shift is finished (relinearized) and rotated on its own; the products of the zero-shift block are summed as with BSGS. Per matvec of dimension d (offset 0): d - 1 rotations, of which only the baby - 1 rotations of x can be hoisted; with lazy\_relin 1 + d - baby relins, without it d (EvalConfig::matvec\_rotations, EvalConfig::matvec\_enc\_relins). At d = 16 (baby = giant = 4), as (hoistable rotations / individual rotations / relins): | form | hoistable rot. | individual rot. | relins (lazy) | relins (not lazy) | |---|---|---|---|---| | default BSGS | 3 | 3 | 4 | 16 | | bsgs = false (this library) | 3 | 12 | 13 | 16 | | textbook diagonal method (not implemented) | 15 | 0 | 1 | 16 | The textbook (Halevi-Shoup) diagonal method rotates x by all d - 1 steps (all hoistable) and sums the products lazily; it needs un-pre-rotated diagonals, i.e. another weight layout and gradient form, so it is not implemented. The all-off configuration has the same counts in both forms (15 rotations, 16 relins at d = 16).\par
2. Bias. Against the strongest naive baseline (the textbook method with hoisting and lazy relinearization), the cost of the bsgs = false rows overstates what BSGS gains: the measured BSGS gain is an upper bound. Conversely, the marginal gains of hoist and lazy\_relin measured inside bsgs = false configurations are lower bounds (only baby - 1 rotations can be hoisted and only the zero-shift block summed lazily: e.g. lazy\_relin saves 3 relins at d = 16 instead of 15).\par
3. hoist changes no OpCounter value. The counter does not distinguish hoisted from individual rotations, so the effect of hoist shows only in wall-clock time, never in key-switch counts.\par
4. bsgs does not touch the gradient's outer product (Dense::weight\_grad\_scaled, Dense::weight\_grad\_reduced): its (baby - 1) + (giant - 1) input rotations are what the pre-rotated weight layout requires, not a BSGS choice; only hoist changes how they run. The interactive gradient multiplies each diagonal with one Evaluator::mul (nothing to sum lazily), so lazy\_relin does not change it either.\par
The smoke rows (insecure N = 2\textasciicircum{}12, high Mac load) check counts and closed forms only; their times are not comparative.\par
\end{minipage}
\end{table*}

Table~\ref{tab:gpu_vs_cpu_ops} contains synchronized operation latencies, including capped-level rows and memory skips. The HEonGPU relinearization reference~\cite{gpuKeys2025} is nominal: its modulus count differs. A linear-in-limb adjustment is a cost-model extrapolation, not an equal-parameter measurement. The published table value is inherited from the archived operator report; the bibliographic metadata was checked for this draft, but a fresh PDF fetch was blocked. The later pool/submission fix has not been retimed.

Table~\ref{tab:rq3_ablation} deliberately leaves the server rows pending. The insecure small-parameter CPU smoke validates counts and checks only; its high-load timings are inadmissible for performance comparisons. The flag named \texttt{bsgs=false} keeps pre-rotated weights rather than implementing the strongest textbook diagonal baseline. Thus its BSGS gain would be an upper bound against that alternative, while the marginal hoisting and lazy-relinearization gains within those rows would be lower bounds. Hoisting changes no operation counter; the gradient outer product retains layout-required rotations independent of the BSGS flag. All four caveats travel with the table.

CPU and GPU training comparisons require the same optimizer, level, layout, and batch. Registered SGD fallback rows are compared only with the corresponding CPU SGD row. A resource skip is unmeasured; a device OOM is a failure. Test-process success cannot stand in for an actually executed production-size parity gate. GPU bootstrap precision, latency, allocation behavior, and peak memory are still unmeasured. HEonGPU's separate bootstrap reference~\cite{heongpuBench} has a different remaining-level budget and main-secret distribution; the fixed plan target is not a normalized speed-ratio result. RQ3's epoch-time question is therefore unanswered.

\section{Related Work, Limitations, and Conclusion}
\label{sec:limits}
\subsection{Positioning}
The earlier diagonal-training design~\cite{mihara2020} motivates rotation-only transposition. We generalize the layout and charge server-side batch maintenance, compact refresh, and optimizer state. ReBoot~\cite{reboot2026} already establishes non-interactive training and local learning; its released implementation is our application baseline, not a source of imported implementation code. HETAL~\cite{hetal2023} studies encrypted transfer learning and efficient matrix products in a different training scope. Our evaluation does not reproduce its workloads or establish a ranking against it.

The cryptographic building blocks are efficient linear transforms and hoisting~\cite{halevi2018}, SSE~\cite{sse2022}, and CKKS~\cite{ckks2017}. SEAL~\cite{seal} provides interoperability checks; OpenFHE~\cite{openfhe2022} supplies the separately executed baseline. HEonGPU~\cite{gpuKeys2025,heongpuBench} is a device-operation reference, not an end-to-end training competitor measured under our protocol. An MIT implementation improves reuse, but licensing and pipeline completeness are not evidence of a cryptographic or learning-performance advance. The approved dependency policy permits MIT, Apache, BSD, and ISC licenses, with one narrow Unicode-3.0 character-table exception for \texttt{unicode-ident} version \texttt{1.0.27}. That component requires its complete Unicode notice; this is not a general Unicode-license exception.

\subsection{Limits and missed goals}
\paragraph{Scope and comparison.}
FULL-1 and FULL-2 supply completed full-setting encrypted runs at two batches with the same initial-weight seed, including their schedules and accumulated error. FULL-2 passes the registered baseline-batch accuracy gate, but misses the common-baseline-bootstrap count target. FULL-1's successful count comparisons are a separate configuration. Full-setting ReBoot CPU wall-clock is unmeasured, and the common-security condition is unmet. Neither these two configurations nor the earlier subsets establish multi-seed encrypted behavior; multi-seed plaintext results are a different evidence class. The observed parameter error does not establish a general error bound. Batch choices trade refresh amortization against learning dynamics. The effective baseline security differs, initialization and bias arithmetic require care, and legacy rows cannot support accuracy claims. ReBoot's in-library layout timing was measured on a different machine and is not a same-CPU baseline; its public-label/client-update path has a weaker confidentiality contract. The official encrypted block-model baselines have not been run: their reported memory demand exceeds the available host budget.

\paragraph{Cost and correctness.}
A key switch at a low level need not cost the same as one at the top level; counts omit memory traffic, cache effects, and key residency. The initial compact-layout assumptions required actual CKKS tests and the first model was optimistic. Replicated copies were once unstable during interactive training; canonical-layout refresh was required to avoid divergence. The earlier paper's published timing total and phase sum disagree, so the continuity demonstration does not silently reconcile them or infer a server-only speedup. The original CPU NTT target and several SEAL-relative key-switch performance targets were missed. The early CUDA NTT spike was not timed against its registered hardware reference; later device-operator measurements do not retroactively supply that measurement. Some SEAL comparisons also have different special-prime/key-switch decompositions. We report these as missed performance targets, not failed numerical gates.

\paragraph{Numerical limits.}
SSE and modulus-band changes were needed because the planned dense-secret overflow range and original prime allocation did not meet the bootstrap precision budget. Unit-scale nonconstant masks were invalid, and exact-scale folding is required to fit the training circuit. The accepted security claim uses specified sieve cost models, not a universal security level; Core-SVP estimates are lower. Dense-secret hybrid search is bounded rather than an exhaustive global optimization. An all-sparse main-key alternative did not meet the selected security criterion and is not used. KAT coverage is limited to the supported SEAL preset, with asymmetric comparison directions; it does not independently validate our bootstrap or larger hybrid preset. Standalone OpenFHE/Lattigo bootstrap precision parity remains absent. Random full-slot CPU precision checks do not prove a negligible failure probability for arbitrary training distributions or activations outside range.

\paragraph{Learning limits.}
The quartic activation's claimed accuracy benefit is unmeasured. The evaluated encrypted trainer omits weight decay and optional softmax/cross-entropy. The separate plaintext paper-setting reconstruction includes weight decay; it is not the faithful encrypted protocol. Bias momentum is not matched to NAG, and faithful bias freezing preserves extra encrypted arithmetic. Fixed-final-epoch Breast Cancer accuracy misses its target; test-selected best epochs are not a remedy. No encrypted Fashion-MNIST or Breast Cancer accuracy is reported. Higher-depth local-block training has plaintext fidelity checks but no completed encrypted comparison. The continuity Iris experiment is interactive and differs in padding, randomization, and normalization; it cannot establish the primary non-interactive claim. Threshold or multi-party FHE, scheme switching, transformer training, and Metal execution are outside the evaluated scope.

\paragraph{Device and implementation limits.}
Full-key GPU training, the matched training parity gate, the RQ3 timing ablation, and GPU bootstrap measurements are pending. Repaired multi-stream latency is measured only for independent products at a capped level; it cannot supply those missing training results. Memory estimates and pool snapshots do not establish a peak. All-off flags do not implement the strongest naive diagonal algorithm. The current training benchmark still uses host bootstrap fallback. Device error handling and numerical replay are useful engineering evidence, but compile-only paths and skipped tests are not physical hardware results. CPU/GPU side-channel resistance, constant-time sampling, secret zeroization, malformed-input hardening, and independent cryptographic audit are outside the current implementation's assurance. Some internal invariants are debug assertions; decode/parameter mismatches can panic. The artifact is research software, not a deployment security claim.

\subsection{Conclusion}
Batched diagonal packing preserves a rotation-only transpose while supporting compact non-interactive optimizer updates. The completed full-setting runs show a larger-batch count advantage and baseline-batch accuracy within tolerance, together with a missed normalized count target at the baseline batch. The security mismatch and unmeasured full-setting baseline time prevent an unconditional superiority claim. The device microbenchmark also shows that stream and pool choices matter, without resolving training performance. The resulting evidence supports a conditional re-evaluation of the layout and identifies the matched-security, baseline-timing, and hardware experiments still needed.

\bibliographystyle{ACM-Reference-Format}
\bibliography{refs}
\appendix
\section{Parameters, Security, and Reproduction}
\label{app:params}
\subsection{Implemented chain and bootstrap contract}
The production bootstrap preset uses $N=2^{16}$, $N/2$ real slots, scale $2^{42}$, and Q-prime bit sizes
\begin{equation*}
 [54,\underbrace{42,\ldots,42}_{16},\underbrace{34,\ldots,34}_{4},
 \underbrace{58,\ldots,58}_{7},\underbrace{56,\ldots,56}_{4}].
\end{equation*}
Four special primes have bit size $60$, with eight hybrid digits. The bands allocate four levels to each linear transform and seven to EvalMod, leaving level $16$ before training unpack. EvalMod uses Chebyshev degree $127$, interval parameter $K=13$, and no double-angle step. These are implemented design constants; the superseded original-plan bands and dense-secret range are not used.

The bootstrap requires level zero and exactly the preset scale before work. Dense-to-sparse SSE is restricted to $q_0P$ under a temporary key of weight $32$; the return key and remaining arithmetic use the dense main key over divisors of $QP$. CPU full-slot tests exceed the registered precision threshold for random real inputs in $[-1,1]$. This does not imply a worst-case failure bound. The GPU implementation has local replay and compilation coverage, but its physical precision and memory feasibility remain unmeasured.

\subsection{What the estimator models}
Our default sampler gives $-1,0,1$ probabilities $1/4,1/2,1/4$. We model the dense secret with a fixed-count sparse-ternary proxy rather than a uniform-ternary distribution; SSE is a separate small-modulus instance. The reported attack input rounds the exposed $\log_2(QP)$ upward and uses $m=N$ samples. Dense and sparse main keys must not be interchanged at the same modulus without re-estimation.

The estimator is pinned to commit \texttt{ae45d740ceb1}; methodology follows Albrecht et al.~\cite{albrecht2015} and the maintained estimator documentation~\cite{latticeEstimator}. The recorded search includes primal/dual and selected hybrid/MITM variants. Dense primal-hybrid optimization did not finish globally; fixed guessing-dimension scans supplement it. Arora--GB and BKW were not included. For the released ReBoot context, an observed sign split from one generated key is an input instance, not a distributional guarantee for all keys. The minima are conditional on the selected search, cost model, and memory assumptions. Table~\ref{tab:security} is not a proof, nor evidence of a practical compromise.

\subsection{Transpose identity and reduction boundary}
For Eq.~\eqref{eq:transpose}, matching row and column indices gives $t\equiv i\pmod {n_o}$ and $t+c\equiv0\pmod {n_i}$. The wide/tall cases in Section~\ref{sec:packing} solve both conditions. Substituting the pre-rotations on the source and destination yields the displayed offset. Because every weight diagonal is periodic with period $L$ in feature coordinates and all offsets are multiples of $B$, samples are not mixed. The identity therefore applies to rectangular padded layers and repeated blocks as well as square matrices. A zero offset is a copy; other offsets require the selected dedicated or composed rotation keys.

Lane reduction is different: offsets not divisible by $B$ mix lanes intentionally, and isolating one lane requires a mask. Ignoring that mask or its replication work changes both the depth and count model. A partial-batch validity mask must also exclude unused lanes from the loss derivative; zero input alone is insufficient when the model has a bias or a linear output loss.

\subsection{KAT and artifact boundaries}
The SEAL compatibility checks compare encoding, serialized ciphertext interoperability, decryption, and supported evaluation paths under the shared preset. Some directions assert exact residues while others only assert a decoded-error bound. Rotation arithmetic intentionally differs in its centered lifting. Neither this KAT nor CPU/GPU bit equality establishes security; the latter only checks implementation agreement for executed cases. Small insecure presets are smoke tests and never enter the security claim.

Tables are generated by \texttt{tools/paper\_tables.py}; field-level provenance is available through its \texttt{--sources} option. The figure generator consumes those validated table cells and writes a source manifest containing paths, fields, and SHA-256 digests. It does not rewrite result JSON. Missing registered runs produce explicit pending rows; existing invalid inputs must fail validation. Timed runs, setup, client diagnostics, and model-only extrapolations stay distinguishable. Running \texttt{scripts/run\_all.sh --tables-only} regenerates the retained tables. This does not execute the complete registered experiment matrix, whose end-to-end runner remains unfinished; a quick pipeline check is not a source of research results.

The anonymous review mode omits author metadata and identifying artifact links. ReBoot and OpenFHE are executed as external baselines; no implementation source is copied into this artifact. The estimator and TeX compiler are isolated development tools, not library dependencies. The library's dependency policy, including the specific character-table exception disclosed in Section~\ref{sec:limits}, is separate from the licenses of those external research tools. Instructions for rebuilding anonymous review, named review, and named preprint variants are provided with the source. Authorship is confirmed for this preprint. Before arXiv submission, the final source, metadata, and server-generated PDF require checking against the retained evidence; unavailable measurements remain unavailable. A later conference submission requires a separate anonymity and venue-rules review; the WAHC~2026 limit is only the provisional formatting constraint~\cite{wahc2026}.

\ifdefined\ArxivDraft
Code and evidence: \url{https://github.com/kenmaro3/ckks-train}. Release \texttt{v1.0.0-preview.1} is an earlier snapshot; this preprint has its own source revision.
\fi

\end{document}
 for placeholders.
\newif\ifshowauthors
\ifdefined\AuthorDraft\showauthorstrue\else\showauthorsfalse\fi
\ifshowauthors
  \documentclass[sigconf,review]{acmart}
\else
  \documentclass[sigconf,review,anonymous]{acmart}
\fi
\setkeys{acmart.cls}{balance=false}
\setcopyright{none}
\settopmatter{printacmref=false}
\acmConference[Draft]{Unsubmitted research draft}{November 2026}{No venue assigned}
\acmDOI{}
\acmISBN{}
\usepackage{amsmath}
\usepackage{graphicx}
\usepackage{booktabs}
\newcommand{\sys}{\textsc{ckks-train}}
\newcommand{\rot}{\operatorname{rot}}
\newcommand{\KS}{\operatorname{KS}}
\newcommand{\ceil}[1]{\left\lceil #1\right\rceil}
\newcommand{\floor}[1]{\left\lfloor #1\right\rfloor}
\title[Diagonal Packing for Encrypted Training]{Revisiting Diagonal Packing for Non-Interactive Encrypted Neural Network Training}
\ifshowauthors
  \author{Author names to be supplied}
  \affiliation{\institution{Affiliation to be supplied}\city{City to be supplied}\country{Country to be supplied}}
  \email{author@example.invalid}
\fi
\begin{document}
\begin{abstract}
Training a neural network under homomorphic encryption requires more than an efficient forward matrix product: transposes, batch reductions, optimizer state, and repeated refreshes can dominate the cost. We revisit diagonal packing for this setting. Starting from the earlier rectangular transpose, we combine a periodic representation using the smaller padded dimension's number of diagonals with feature-major SIMD mini-batch lanes and baby-step giant-step evaluation. Compact optimizer state amortizes bootstrapping across diagonals. We implement the resulting training path in a Rust CKKS stack with instrumented CPU and CUDA arithmetic. We compare key-switch accounting conventions, bootstrap placement, and measured operator latencies, and explicitly separate layout effects from bootstrap implementations. The available encrypted subset experiments support a conditional key-switch advantage for a larger batch, but do not settle the registered accuracy and equal-security comparison. Full-setting encrypted results, GPU training ablations, and physical GPU bootstrap measurements remain pending. We report these gaps together with negative results rather than interpreting implementation completeness as research success. This interim draft establishes the algorithms and reproducible evidence needed to finalize that comparison.
\end{abstract}
\begin{CCSXML}
<ccs2012><concept><concept_id>10002978.10002979.10002980</concept_id><concept_desc>Security and privacy~Cryptography</concept_desc><concept_significance>500</concept_significance></concept><concept><concept_id>10010147.10010257.10010293.10010294</concept_id><concept_desc>Computing methodologies~Neural networks</concept_desc><concept_significance>300</concept_significance></concept></ccs2012>
\end{CCSXML}
\ccsdesc[500]{Security and privacy~Cryptography}
\ccsdesc[300]{Computing methodologies~Neural networks}
\keywords{homomorphic encryption, CKKS, encrypted training, diagonal packing, bootstrapping}
\maketitle
\section{Introduction}
\label{sec:intro}
Homomorphic encryption allows a service to evaluate a public computation without learning the encrypted inputs. Training is particularly demanding: the model changes at each step, gradients require transposed products, and a long optimization trajectory repeatedly exhausts the available multiplicative depth. A useful packing scheme must therefore account for the entire update, not just inference.

Mihara et al.~\cite{mihara2020} showed that diagonal-packed weights support transposition through rotations rather than ciphertext multiplications. Their interactive training design used a client to restore ciphertexts between updates. ReBoot~\cite{reboot2026} subsequently demonstrated non-interactive CKKS neural-network training, combining packed arithmetic, local learning signals, and bootstrapping. Non-interactivity itself is therefore not a novelty claim of this work. Our question is whether diagonal packing remains competitive when batch reductions, server-side layout maintenance, optimizer state, and bootstrapping are all charged.

We develop \sys{}, an MIT-licensed Rust implementation, and study three questions. \textbf{RQ1} asks how many key switches a complete training update requires compared with a ReBoot-style layout and the released ReBoot baseline. \textbf{RQ2} asks how bootstrap placement interacts with depth and compact state. \textbf{RQ3} asks what CUDA execution and its constituent optimizations contribute to measured training cost. A secondary continuity experiment revisits the original Iris setting; it is an interactive demonstration, not evidence of non-interactive training performance.

Our contributions are: (i) a rectangular, batched diagonal layout with a rotation-only transpose and explicit reduction costs; (ii) a compact-state non-interactive update whose operation counts can be checked against structural predictions; and (iii) an implementation and evaluation protocol that separates cryptographic assumptions, counter-derived comparisons, measured timing, and missing runs. Hoisting, baby-step giant-step (BSGS), sparse-secret encapsulation (SSE), and bootstrapping are established techniques; the contribution lies in their training-specific composition and accounting, not in inventing these primitives.

\paragraph{Status of the evidence.}
This is an interim full draft. The generated tables retain pending rows for the registered full-setting encrypted experiment and GPU training ablation. The completed encrypted subset runs use a different input dimension and cannot replace those rows. The released baseline and our preset also have different secret distributions and security estimates. We consequently do not claim that the registered equal-security, accuracy-matched research goal has been met. Section~\ref{sec:limits} identifies what would change that conclusion.

\section{Preliminaries and Threat Model}
\label{sec:prelim}
\paragraph{Arithmetic.}
CKKS~\cite{ckks2017} encodes approximate vectors in
$R_Q=\mathbb{Z}_Q[X]/(X^N+1)$, with up to $N/2$ complex slots. We use real slot values. A ciphertext at level $\ell$ uses $q_0,\ldots,q_\ell$ and a scale $\Delta$. Multiplication followed by rescaling consumes a modulus prime. Rotations permute slots through a Galois automorphism and a key switch. Hybrid RNS key switching uses special primes $P$ and a digit decomposition; its cost depends on the active limbs, digit sizes, and memory traffic.

We count rotations (including conjugations), relinearizations, and explicit SSE switches:
\begin{equation}
  \KS = n_{\mathrm{rot}}+n_{\mathrm{relin}}+n_{\mathrm{SSE}}.
  \label{eq:ks}
\end{equation}
A ciphertext tensor product is not automatically a key switch: lazy relinearization may combine several products before one switch. Bootstrap events are recorded separately and their constituent operations are already in Eq.~\eqref{eq:ks}; adding a bootstrap-equivalent charge again would double count. Reported counts cover the whole network and are averaged over processed examples, not over layers. A count is a workload metric, not a machine-independent latency estimate.

\paragraph{Refreshes.}
An interactive refresh decrypts and re-encrypts at the key owner. A bootstrap instead evaluates a refresh circuit using public evaluation keys. Our non-interactive path uses SSE~\cite{sse2022}, centered modulus raising, coefficient-to-slot and slot-to-coefficient transforms, and an approximate modular-reduction polynomial. Public transforms use collapsed stages and BSGS; exact scale tracking makes them composable with training.

\paragraph{Protected objects and public leakage.}
We consider a semi-honest evaluation server and a single trusted key owner, under the usual RLWE and evaluation-key assumptions of the chosen CKKS construction. No new security reduction is claimed. Training examples, labels, model weights, intermediate activations, gradients, and optimizer state are encrypted. The architecture, dimensions, batch schedule, hyperparameters, parameters, evaluation keys, and algorithmic control flow are public. The key owner initializes and encrypts the data and initial model, and decrypts the resulting model. Known initialization distributions are not themselves secret. This model does not hide dimensions, resource use, access schedules, or timing, and it does not provide malicious-server verification, differential privacy, multi-party key custody, or protection against inference from a legitimately decrypted model.

\paragraph{Diagnostics are not a server capability.}
The benchmark process can also instantiate a secret-key holder, decrypt checkpoints and predictions, compare against a plaintext twin, and report accuracy or activation ranges. Those are client-side diagnostics for a trusted experiment harness; they must not be exposed as server logs in a deployment. The generic non-interactive trainer has no secret key and makes public depth-based refresh decisions. Disabling diagnostic decryption changes neither the encrypted arithmetic nor the public schedule. The interactive path explicitly involves the client during training and is labeled separately. The in-library ReBoot-layout microbenchmark uses plaintext labels and client-side optimizer updates: it supports arithmetic/count comparisons, not the full confidentiality claim or a faithful non-interactive baseline.

\section{Rotation-Only Transposition for Mini-Batch Training}
\label{sec:packing}
\subsection{A common layout for activations and weights}
Let $B$ be a power-of-two lane count and $d$ a power-of-two feature capacity such that $dB$ divides the slot count. Slot $jB+b$ stores feature $j$ of sample $b$. Blocks repeat when $dB$ is smaller than the slot count. A layer has positive integer input and output dimensions $d_i,d_o$. Define the padded dimensions $n_i=2^{\ceil{\log_2 d_i}}$ and $n_o=2^{\ceil{\log_2 d_o}}$, and require $n_i,n_o\le d$. A vector of padded dimension $n\in\{n_i,n_o\}$ is canonical when its feature sequence is $n$-periodic in every lane; unused logical features are zero. This periodicity eliminates a separate replication step at tall layers.

Zero-pad the original $d_o\times d_i$ weight matrix to $n_o\times n_i$. Let $m=\min(n_o,n_i)$ and $L=\max(n_o,n_i)$. Both divide $d$, so the $L$-periodic diagonals, rotations modulo $LB$, and integer logarithmic folds below are well-defined. Store the generalized diagonals
\begin{equation}
 D_c[j]=W[j\bmod n_o,(j+c)\bmod n_i],\quad 0\le c<m.
 \label{eq:diagonal}
\end{equation}
Weights are shared across all lanes. With BSGS split $m=gu$ and $c=ga+b$, the stored diagonal is
$P_c=\rot(D_c,-gaB)$. The encrypted product is
\begin{equation}
 y=\sum_{a<u}\rot\!\left(\sum_{b<g}P_{ga+b}\odot\rot(x,bB),gaB\right).
 \label{eq:bsgs}
\end{equation}
Here a positive rotation reads slot $s+k$ into output slot $s$; $\odot$ denotes a slot-wise product. For a wide layer, $\log_2(n_i/n_o)$ rotations and additions fold partial outputs together. For a tall layer, the canonical input already has the required replication. Neither case adds a mask to the fold.

\subsection{Transpose without products or masks}
Rectangular diagonal transposition without ciphertext multiplications is inherited from Mihara et al.~\cite{mihara2020}. Their general method uses rotations and additions; their implementation also reports OpenMP batch processing. Here we compose that idea with the periodic $m$-diagonal representation, feature-major SIMD lanes, and BSGS pre-rotations.

The transpose diagonal $E_i[j]=W[(j+i)\bmod n_o,j\bmod n_i]$ equals a rotation of one $D_c$. Choose $c=(-i)\bmod m$, and set $t=-c$ for a wide layer or $t=i$ for a tall layer. In stored BSGS form,
\begin{equation}
 Q_i=\rot(P_c,\,[t-g\floor{i/g}+g\floor{c/g}]B).
 \label{eq:transpose}
\end{equation}
The rotation amount is taken modulo $LB$. Thus transposition uses neither a ciphertext product nor a mask and consumes no multiplicative level. It is \emph{not} free: each nonzero rotation requires a key switch. Dedicated keys minimize rotation count; composing existing baby/giant keys saves key memory at the expense of extra switches. Our training runs use key reuse, and the extra work is included in the counters.

\subsection{Gradients, reductions, and a symbolic cost model}
For layer gradient $\delta$ and input $a$, a pre-rotated weight-gradient diagonal is
\begin{equation}
 \nabla P_{ga+b}
 =\operatorname{batchsum}\bigl(\rot(\delta,-gaB)\odot\rot(a,bB)\bigr).
 \label{eq:gradient}
\end{equation}
Batch summation first reduces lanes by powers-of-two rotations. A lane-zero mask isolates the reduced value; inverse-direction rotations replicate it when replicated gradients are required. The mask has nonintegral polynomial coefficients in the CKKS embedding, so encoding it at unit scale is not an exact, zero-depth shortcut. It consumes a level. Constants for the learning rate and batch normalization are folded into a necessary plaintext multiplication at its exact target scale.

\begin{table}[t]
\centering\small
\caption{Symbolic layer costs for the implemented layout, before any wide-layer fold. $g,u$ are the BSGS split, $m=gu$. Gradient costs show the replicated interactive form; compact non-interactive state changes the final replication.}
\label{tab:symbolic}
\begin{tabular}{lcc}
\toprule
Operation & Rotations & Relins \\
\midrule
BSGS encrypted matvec & $g+u-2$ & $u$ \\
Dedicated transpose & $m-1$ & $0$ \\
Replicated weight gradient & $g+u-2+2m\log_2 B$ & $m$ \\
\bottomrule
\end{tabular}
\end{table}
Hoisting reuses decomposition for rotations of the same input~\cite{halevi2018}. It changes wall-clock work but not Eq.~\eqref{eq:ks}. Lazy relinearization applies within a matvec giant-step sum. It cannot remove the separate relinearizations needed before gradient-diagonal rotations. These distinctions matter when interpreting ablation results.

\paragraph{Extensions relative to the interactive design.}
The first two extensions studied here are the periodic $\min(n_o,n_i)$-diagonal representation with feature-major SIMD mini-batch lanes, and its composition with hoisting and BSGS. General rectangular transposition and batch processing themselves are inherited capabilities, not novelty claims. The remaining four extensions---bootstrap placement, polynomial activations, momentum, and loss conventions---are specified in Section~\ref{sec:noninteractive}. None removes the cost of maintaining canonical server-side state; the initial interactive cost model omitted some replication work and was therefore too optimistic.

\section{Non-Interactive Training}
\label{sec:noninteractive}
\subsection{Compact state and the complete update}
Refreshing every replicated diagonal separately wastes bootstrap slots. We instead reduce each gradient to one lane per feature and combine masking, scaling, and lane shifting in a plaintext product. Selected lanes from several diagonals form compact weight and momentum containers. A public compact-layout plan fixes their offsets and the required rotations; unpacking restores the periodic diagonals for the next forward pass. Bias slots can share unused compact positions. Masked packing and unpacking both have a depth cost and are included in the closed-form predictor.

The server performs encrypted forward propagation, forms the loss gradient from encrypted labels, propagates through the transposed weights, and accumulates gradients over micro-batches before one optimizer update. Increasing the effective batch beyond the physical lane count therefore amortizes optimizer refreshes without pretending the examples fit in one ciphertext. It also changes optimization behavior and cannot be treated as a cost-only substitution.

The original dense-secret bootstrap constants did not satisfy the desired precision and depth budget. The implemented preset instead uses a dense ternary main key and a temporary sparse key during modulus raising. The SSE pair is returned to the dense main key before the remaining arithmetic. Collapsed transforms and a level-optimal Chebyshev evaluator fit the selected modulus bands. This is an engineering parameter choice, not a new bootstrap construction or an inherited failure-probability theorem. Appendix~\ref{app:params} specifies the chain and the exposed key instances.

\subsection{Three refresh schedules (extension 3)}
\texttt{AfterUpdate} (AU) bootstraps every compact weight and velocity container after each update. \texttt{WhenNeeded} (WN), an additional schedule beyond the original plan, does so only if the next update would exceed the remaining depth. The decision uses the next batch's fullness: padded lanes at a linear output need a validity mask, so a partial batch can consume an extra level. An unknown next shape is conservatively treated as partial.

\texttt{AfterActivation} (AA) refreshes each hidden activation and retains the WN rule for compact weights. For the shallow eMLP-1 workload, activation refresh does not reduce the frequency of weight refresh beyond WN; it adds a bootstrap per hidden activation and micro-batch. Its motivation is enabling deeper circuits that otherwise exceed the budget, not an automatic speedup. Activation magnitudes must fit the supported bootstrap range. The observed AA subset run exceeded the conservative range in some slots and had a larger activation error; it was not clipped, and does not certify arbitrary activation distributions.

For a sequence of identical updates consuming $D$ levels with $L_a$ usable levels after unpacking, $\floor{L_a/D}$ is only an initial scheduling bound. Unpack costs between updates, partial batches, and the actual state level must also enter the scheduler. We check the structural predictor against every measured update rather than replacing this stateful schedule with a fixed average.

\subsection{Activation, optimizer, and loss (extensions 4--6)}
\paragraph{Polynomial activation.}
We retain the earlier squared activation for continuity and implement $z^2+z$ for the ReBoot comparison. A polynomial quartic option is available for depth experiments. Its additional depth and available tests do not establish a measured accuracy improvement; that planned claim remains unsupported.

\paragraph{Momentum.}
Both SGD and Nesterov-style momentum are supported. Writing $h=\eta(1+\mu)G/n$ and scaled velocity $U=\eta(1+\mu)V$ gives
\begin{align}
 U'&=\mu U+h,\nonumber\\
 W'&=W-h-\frac{\mu^2}{1+\mu}U.
 \label{eq:nag}
\end{align}
Constants are folded into exact-scale products already needed by the gradient path. A changing learning rate requires rescaling the stored velocity representation, and its lane shifts are charged. Biases use SGD; the faithful ReBoot configuration freezes their encrypted zero values while retaining the implementation's bias arithmetic. Weight decay is not implemented; the captured faithful setting uses zero decay.

\paragraph{Loss.}
A squared-error loss yields an affine output gradient; MSE normalization and ReBoot's sum reduction are matched through the learning-rate convention. The faithful configuration uses the released initialization, sample order, output padding, and schedule, not just the paper's nominal hyperparameter names. Softmax/cross-entropy approximation is optional in the plan and is not implemented or evaluated here. Block MLPs follow local classifier losses and stop gradients across blocks; they must not be described as an ordinary deeper back-propagation network.

\section{GPU Implementation}
\label{sec:gpu}
\sys{} separates RNS arithmetic, CKKS evaluation, neural-network operations, CUDA execution, and benchmark orchestration. Generic \texttt{HeOps} and \texttt{BootstrapOps} interfaces let the same trainer use CPU or device ciphertexts. The evaluator checks parameter identity, levels, scales, key availability, and device ownership before work. Counters commit after success; the CUDA bootstrap also waits for final device synchronization before merging its internal counters.

The CUDA backend uses device NTTs, basis conversion, pointwise arithmetic, hybrid key switching, and BSGS matvecs. Small public launch data use pinned transfer buffers. Reusing a hoisted decomposition avoids repeating the expensive preparation for related rotations. Device/CPU bit parity is meaningful only for cases actually executed; memory-skipped test functions can still report process success.

\paragraph{Streams and memory.}
Independent products may be issued to several streams, but a training step's dependent matvecs are executed sequentially. The measured pre-fix multi-stream benchmark was slower than sequential execution (Table~\ref{tab:gpu_vs_cpu_ops}). Calling-thread submission and pool handling were subsequently revised; their performance remeasurement is pending. An effective-bandwidth estimate for addition does not prove matvec occupancy or full-device saturation. That explanation remains a hypothesis until profiled.

The shared device has limited available memory. Training can upload keys capped to an explicitly selected level; this preserves the computation at that level, but does not measure the top-level workload. Allocation preflight uses estimated resident/key/workspace sizes. Pool trims occur after large setup stages, at epoch boundaries, and between ablation configurations. Estimated need and pool checkpoints are not a measured memory peak, and a successful preflight does not guarantee later allocations under contention.

\paragraph{Bootstrap boundary.}
The CUDA bootstrap implementation keeps ciphertext arithmetic on the device through SSE, modulus raising, transforms, and EvalMod. A narrow shared polynomial executor preserves the CPU operation order and rounding. Public integer constants are still materialized on the host and uploaded; future device timing must include that overhead. Local CPU replay and CUDA compilation support code readiness, not physical-GPU precision, latency, or memory feasibility. Those hardware measurements are pending.

The training benchmark still uses a CPU bootstrap adapter. Level operations run on the GPU; non-interactive refresh downloads to the CPU and uploads the result. A compiled device bootstrap API has not changed that benchmark. Its timing must therefore separate device work, transfers, and CPU refresh. Likewise interactive GPU training uses a client refresh. Neither path justifies an end-to-end all-GPU training claim.

\section{Evaluation}
\label{sec:evaluation}
\subsection{Protocol and evidence classes}
We distinguish a released baseline execution, an in-library layout microbenchmark, and a cost model. The first supplies observed wall-clock and bootstrap events; the second isolates arithmetic layout counts under our evaluator; the third supplies uninstrumented baseline quantities. A counter-derived normalization is not a separate execution. Historical generated tables retain the working name \emph{hekit} for \sys{} and its crate names.

The headline uses MNIST reduced to $196$ inputs and the released full training setting. FULL-1 has completed at $B=64$ with one initial-weight seed; FULL-2 at the released baseline batch remains pending. FULL-1 is partial full-setting evidence, not a substitute for the FULL-2 accuracy gate. The completed faithful encrypted subset uses $784$ inputs, the baseline's sample indices and initial weights, and its training conventions. Output width is padded to $16$, including the unused classes in the released loss; replacing this with a width-$10$ loss is not assumed equivalent. Earlier ``legacy'' runs use a different learning rate, bias policy, and initialization. Their counters and timing remain useful, but their accuracy is not comparable to the released baseline. The last subset batch is partial, affecting depth, refresh scheduling, and the effective sum-versus-mean learning-rate conversion.

The paired server measurements use the same CPU and thread cap, with different CPU parallel runtimes. Jobs are not core-pinned, and some rows overlap unrelated shared-host load. Client work and setup scopes also differ: baseline step time includes batch encryption and prediction decryption, whereas our update time excludes those separately reported client operations; both exclude key setup. Thus even a same-host ratio is not a controlled architectural speedup. We do not pool CPU and GPU samples. Encrypted rows have a single seed unless the table states otherwise; within-run iterations and Criterion samples are not independent seeds.

\subsection{RQ1: key switches and wall-clock}
% generated by tools/paper_tables.py from results/; do not edit
\begin{table*}[t]
\centering
\scriptsize
\caption{RQ1: key switches per sample, four ways (ratio hekit / ReBoot at B = 48)}
\label{tab:rq1_ks}
\resizebox{\linewidth}{!}{%
\begin{tabular}{lrrrrrrrrrrrrr}
\toprule
Run & Input & Placement & B & Config & Samples & hekit KS/sample & Boots & ReBoot non-boot & (a) & (b) & (c) & (d) & Ways $\leq$ 0.70 \\
\midrule
\multicolumn{14}{l}{\textit{196 inputs, ReBoot full setting, 4800 samples x 10 epochs, seed 0 (headline)}} \\
FULL-1 & 196 & WN & 64 & faithful & 48000 & 18.272 & 1500 & 37 (model) & 0.232 & 0.387 & 0.395 & 0.400 & abcd \\
FULL-2 & 196 & WN & 48 in 64 & faithful & 48000 & 28.990 & 4000 & 37 (model) & 0.368 & 0.781 & 0.627 & 0.534 & acd \\
\midrule
\multicolumn{14}{l}{\textit{784 inputs, ReBoot encrypted-script subset, 1 epoch (supplementary)}} \\
F1 & 784 & WN & 32 & faithful & 1008 & 33.488 & 68 & 35.0625 (Lazy, measured) & 0.435 & 0.779 & 0.756 & 0.742 & a \\
F2 & 784 & WN & 64 & faithful & 1008 & 17.885 & 36 & 35.0625 (Lazy, measured) & 0.233 & 0.414 & 0.404 & 0.397 & abcd \\
R-41-3 & 784 & WN & 32 & legacy (R-M3-12) & 1024 & 32.531 & 64 & 35.0625 (Lazy, measured) & 0.423 & 0.741 & 0.734 & 0.730 & a \\
R-41-4 & 784 & WN & 64 & legacy (R-M3-12) & 1024 & 17.172 & 32 & 35.0625 (Lazy, measured) & 0.223 & 0.382 & 0.388 & 0.391 & abcd \\
R-41-1 & 784 & AU & 32 & legacy (R-M3-12) & 1024 & 39.469 & 128 & 35.0625 (Lazy, measured) & 0.513 & 1.149 & 0.891 & 0.730 & a \\
R-41-2 & 784 & AU & 64 & legacy (R-M3-12) & 1024 & 20.641 & 64 & 35.0625 (Lazy, measured) & 0.268 & 0.586 & 0.466 & 0.391 & abcd \\
R-41-5 & 784 & AA & 32 & legacy (R-M3-12) & 1024 & 36.000 & 96 & 35.0625 (Lazy, measured) & 0.468 & 0.945 & 0.812 & 0.730 & a \\
\midrule
\multicolumn{14}{l}{\textit{Interactive hekit, client refresh, no bootstraps (separate row, not way (d))}} \\
interactive & 784 & --- & 32 & legacy (M2) & 1024 & 19.781 & 0 & 35.0625 (Lazy, measured) & --- & --- & --- & --- & --- \\
\bottomrule
\end{tabular}}
\par\smallskip\begin{minipage}{\linewidth}\scriptsize
Independent interactive reference: ratio 0.564 against ReBoot Lazy; client refresh, not NI-derived way (d) and not a four-way pass result.\par
KS = rotations + relinearizations + SSE key switches (OpCounter). hekit bootstrap = 111 KS (measured, setup.boot\_ks). ReBoot bootstrap = 502 KS and 196-input non-boot = 37 KS/sample are M0 cost-model values (OpenFHE does not count key switches); ReBoot bootstraps 4 ciphertexts per update at B = 48 (measured on the official subset run: 84 bootstraps / 21 steps).\par
(a) own bootstrap costs; (b) both 502; (c) both 111; (d) bootstraps excluded on both sides. (b)-(d) are derived from the measured counters, not separate executions. Criterion: ratio $\leq$ 0.70 (research plan \S{}4.3.2).\par
784 rows are the supplementary subset rows (R-M3-13): ReBoot non-boot = Lazy 35.06, the in-library measurement of ReBoot's layout (Task 33, baseline c). Legacy rows use the pre-R-M3-12 config; their counts equal the faithful config but their accuracy is not comparable.\par
\end{minipage}
\end{table*}

% generated by tools/paper_tables.py from results/; do not edit
\begin{table*}[t]
\centering
\scriptsize
\caption{Wall-clock on the same CPU (quadruped-gpu, i9-13900KF, 12 threads)}
\label{tab:wallclock_same_cpu}
\resizebox{\linewidth}{!}{%
\begin{tabular}{lrrrrrrrrr}
\toprule
Run & Input & B & Samples & Train s & s/sample & vs ReBoot & Boot s/sample & Non-boot s/sample & Threads \\
\midrule
\multicolumn{10}{l}{\textit{196 inputs, full setting (headline)}} \\
FULL-1 & 196 & 64 & 48000 & pending (R-41-FULL-1) &  & ReBoot 196 full: not measured &  &  &  \\
FULL-2 & 196 & 48 in 64 & 48000 & pending (R-41-FULL-2) &  & ReBoot 196 full: not measured &  &  &  \\
\midrule
\multicolumn{10}{l}{\textit{784 inputs, ReBoot encrypted-script subset}} \\
ReBoot official (OpenFHE) & 784 & 48 & 1008 & 3599.9 & 3.571 & 1 & 2.585 & 0.986 & OpenMP 12 \\
F1 & 784 & 32 & 1008 & 3016.7 & 2.993 & 0.838 & 0.989 & 2.003 & rayon 12 \\
F2 & 784 & 64 & 1008 & 1607.9 & 1.595 & 0.447 & 0.522 & 1.073 & rayon 12 \\
R-41-3 & 784 & 32 & 1024 & 2773.6 & 2.709 & 0.758 & 0.871 & 1.838 & rayon 12 \\
R-41-4 & 784 & 64 & 1024 & 1434.8 & 1.401 & 0.392 & 0.435 & 0.967 & rayon 12 \\
R-41-1 & 784 & 32 & 1024 & 4339.1 & 4.237 & 1.187 & 1.736 & 2.501 & rayon 12 \\
R-41-2 & 784 & 64 & 1024 & 2228.9 & 2.177 & 0.609 & 0.867 & 1.310 & rayon 12 \\
R-41-5 & 784 & 32 & 1024 & 3606.6 & 3.522 & 0.986 & 1.316 & 2.206 & rayon 12 \\
\bottomrule
\end{tabular}}
\par\smallskip\begin{minipage}{\linewidth}\scriptsize
ReBoot time = sum of its own step\_time (encrypt batch, forward/backward, decrypt prediction, weight bootstrap; setup excluded). hekit time = time\_s.train\_total (all updates including bootstraps; key generation and bootstrapper setup excluded). Present rows fail closed unless the binary's provenance names the same host and thread count as the ReBoot run.\par
Shared machine: other users' jobs (SLAM, \textasciitilde{}4 cores) ran during some rows; see docs/m3-report.md \S{}3.2. ReBoot's full 196-input setting was not run (estimated 47.6 h), so the headline has no same-CPU wall-clock ratio.\par
\end{minipage}
\end{table*}

Let $K$ be measured total key switches, $H$ bootstrap events, and $S$ examples processed. Let $c_h$ be our measured per-bootstrap cost and $c_r$ the baseline bootstrap model; let $R$ denote the baseline non-bootstrap count per example, $r$ its refreshes per update, and $B_r$ its own batch. The four reported ratios are
\begin{align}
 \rho_a&=\frac{K/S}{R+rc_r/B_r},&
 \rho_b&=\frac{(K-c_hH+c_rH)/S}{R+rc_r/B_r},\nonumber\\
 \rho_c&=\frac{K/S}{R+rc_h/B_r},&
 \rho_d&=\frac{(K-c_hH)/S}{R}.
 \label{eq:ratios}
\end{align}
Convention (a) uses each side's own bootstrap cost. Convention (b) charges both the baseline cost; (c) charges both our cost; (d) excludes bootstrap arithmetic on both sides. The interactive comparison is separate from (d): omitting bootstrap arithmetic does not remove non-interactive compact packing and unpacking. ReBoot remains at its own batch in the denominator; a hypothetical same-batch denominator would answer a different question.

The baseline's bootstrap cost remains a model because the measured OpenFHE execution does not expose key-switch counters. Its observed bootstrap frequency is used, not replaced with our schedule. The subset non-bootstrap reference is the in-library Lazy layout measurement; the full-setting reference is still a model. These boundaries are visible in Table~\ref{tab:rq1_ks}.

% Generated by tools/figs.py; do not edit. See rq1_ratios.sources.json.
\begin{figure}[t]
\centering
\setlength{\unitlength}{1pt}
\begin{picture}(240,143)
\put(28,32){\line(1,0){204}}
\put(28,32){\line(0,1){82}}
\put(5,124){\makebox(225,10){\footnotesize KS ratio to ReBoot (lower is better)}}
\put(25,32.00){\line(1,0){3}}
\put(0,29.00){\makebox(22,7)[r]{\tiny 0}}
\put(25,52.00){\line(1,0){3}}
\put(0,49.00){\makebox(22,7)[r]{\tiny 0.25}}
\put(25,72.00){\line(1,0){3}}
\put(0,69.00){\makebox(22,7)[r]{\tiny 0.5}}
\put(25,92.00){\line(1,0){3}}
\put(0,89.00){\makebox(22,7)[r]{\tiny 0.75}}
\put(25,112.00){\line(1,0){3}}
\put(0,109.00){\makebox(22,7)[r]{\tiny 1}}
\multiput(28,88.00)(5,0){41}{\line(1,0){2}}
\put(176,91.00){\makebox(55,7)[r]{\tiny threshold 0.7}}
\put(42,32){\color[HTML]{315C8C}\rule{14pt}{34.800pt}}
\put(38,68.800){\makebox(22,7){\tiny 0.435}}
\put(42,21){\makebox(14,7){\scriptsize (a)}}
\put(61,32){\color[HTML]{487F9C}\rule{14pt}{62.320pt}}
\put(57,96.320){\makebox(22,7){\tiny 0.779}}
\put(61,21){\makebox(14,7){\scriptsize (b)}}
\put(80,32){\color[HTML]{6C9B8B}\rule{14pt}{60.480pt}}
\put(76,94.480){\makebox(22,7){\tiny 0.756}}
\put(80,21){\makebox(14,7){\scriptsize (c)}}
\put(99,32){\color[HTML]{A2A76B}\rule{14pt}{59.360pt}}
\put(95,93.360){\makebox(22,7){\tiny 0.742}}
\put(99,21){\makebox(14,7){\scriptsize (d)}}
\put(35,6){\makebox(90,10){\footnotesize F1: B = 32}}
\put(145,32){\color[HTML]{315C8C}\rule{14pt}{18.640pt}}
\put(141,52.640){\makebox(22,7){\tiny 0.233}}
\put(145,21){\makebox(14,7){\scriptsize (a)}}
\put(164,32){\color[HTML]{487F9C}\rule{14pt}{33.120pt}}
\put(160,67.120){\makebox(22,7){\tiny 0.414}}
\put(164,21){\makebox(14,7){\scriptsize (b)}}
\put(183,32){\color[HTML]{6C9B8B}\rule{14pt}{32.320pt}}
\put(179,66.320){\makebox(22,7){\tiny 0.404}}
\put(183,21){\makebox(14,7){\scriptsize (c)}}
\put(202,32){\color[HTML]{A2A76B}\rule{14pt}{31.760pt}}
\put(198,65.760){\makebox(22,7){\tiny 0.397}}
\put(202,21){\makebox(14,7){\scriptsize (d)}}
\put(138,6){\makebox(90,10){\footnotesize F2: B = 64}}
\end{picture}
\caption{Supplementary subset RQ1 ratios (one seed per row), generated from the same cells as Table~\ref{tab:rq1_ks}. Bars (a)--(d) denote the four accounting conventions in Section~\ref{sec:evaluation}. These are not equal-security or full-setting results.}
\Description{Grouped bars for the two faithful subset batches under four key-switch accounting conventions. The larger batch is below the threshold in all four; the smaller batch only under own bootstrap costs. One seed per row; not an equal-security comparison.}
\label{fig:rq1-ratios}
\end{figure}

The faithful larger-batch subset row meets the registered count threshold under every reported convention, while the smaller batch meets it only under the own-bootstrap convention. The generated figure makes that dependence visible. The corresponding larger-batch timing is also favorable on the shared CPU (Table~\ref{tab:wallclock_same_cpu}), but its subset accuracy drops relative to the baseline, including in the plaintext twin. Consequently this observation is a tradeoff, not a successful accuracy-matched RQ1 result. FULL-1 also meets the count threshold under all four conventions. Its completed training time has no measured full-setting ReBoot CPU denominator. The pending FULL-2 accuracy gate, missing baseline timing, and security mismatch below prevent a final verdict.

\subsection{RQ2: placement and depth}
The AU, WN, and AA subset rows in Tables~\ref{tab:rq1_ks} and~\ref{tab:wallclock_same_cpu} provide the placement comparison. For matched legacy eMLP-1 settings, WN reduces bootstrap events and total time relative to AU. AA leaves the same lazy weight-refresh frequency and adds activation refreshes, so it does not improve this shallow workload. Its benefit is a feasible depth budget for some deeper polynomial circuits, checked separately as a structural test rather than an observed training-accuracy gain.

The corrected structural predictor equals the measured counters at every completed update. This validates accounting for the executed paths; it does not validate the baseline bootstrap model or a convergence claim. The initial cost model missed replication, key-reuse, bias, and schedule-related terms, and its optimistic interactive estimate is superseded by that explicit accounting. FULL-1 extends exact schedule agreement to its complete registered update sequence and records accumulated weight and bias error against its same-batch plaintext twin. That observed error is larger than in the earlier subset checks; close final accuracy does not imply a uniform small parameter error. FULL-2 and guarantees across seeds or arbitrary trajectories remain unmeasured.

\subsection{Accuracy, local blocks, and security}
% generated by tools/paper_tables.py from results/; do not edit
\begin{table*}[t]
\centering
\scriptsize
\caption{Test accuracy against ReBoot (R-M3-15)}
\label{tab:accuracy_vs_reboot}
\resizebox{\linewidth}{!}{%
\begin{tabular}{lrrrrrrr}
\toprule
Row & Setting & Seeds & Test acc (\%) & Twin (\%) & Reference (\%) & $\Delta$ (pt) & Gate \\
\midrule
\multicolumn{8}{l}{\textit{Full setting: 196 inputs, 4800 samples x 10 epochs, B = 48 (the $\pm$0.5 pt gate)}} \\
ReBoot released code, plaintext (recomputed twin) & full & 10 & 90.65 $\pm$ 0.27 (n = 10) & --- & --- & --- & reference \\
hekit plaintext twin (ReBoot weights, order, lr) & full & 10 & 90.65 $\pm$ 0.27 (n = 10) & --- & 90.65 $\pm$ 0.27 (n = 10) & +0.00 & twin check \\
hekit encrypted, B = 48 in 64 lanes, WN & full & 1 (seed 0) & pending (R-41-FULL-2) & pending (R-41-FULL-2) & 90.41 &  & pending \\
hekit encrypted, B = 64, WN & full & 1 (seed 0) & 90.36 & 90.37 & 90.41 (seed 0) & -0.05 & not the gate (B $\neq$ 48) \\
\midrule
\multicolumn{8}{l}{\textit{eMLP-2 / eMLP-3 block MLPs, full setting, seed 0}} \\
hekit plaintext twin emlp2 (blocks 64) & full & 1 (seed 0) & 91.56 & --- & 91.56 (ReBoot twin) & +0.00 & twin check \\
hekit plaintext twin emlp3 (blocks 128,64) & full & 1 (seed 0) & 91.75 & --- & 91.75 (ReBoot twin) & +0.00 & twin check \\
\midrule
\multicolumn{8}{l}{\textit{784 inputs, ReBoot encrypted-script subset, 1008 samples, 1 epoch (not a gate, R-M3-15)}} \\
ReBoot official encrypted (OpenFHE) & subset, B = 48 & 1 (seed 0) & 75.06 & --- & --- & --- & reference \\
hekit encrypted WN, B = 32 & subset & 1 (seed 0) & 75.77 & 75.77 & 75.06 & +0.71 & not a gate \\
hekit encrypted WN, B = 64 & subset & 1 (seed 0) & 73.58 & 73.59 & 75.06 & -1.48 & not a gate \\
hekit encrypted eMLP-2, B = 48 (196) & subset 1008 & 1 & pending (R-52-1) & pending (R-52-1) &  &  & twin check \\
hekit encrypted eMLP-3, B = 48 (196) & subset 1008 & 1 & pending (R-52-2) & pending (R-52-2) &  &  & twin check \\
\midrule
\multicolumn{8}{l}{\textit{Reference only: the paper value and the released code at the paper setting (R-REB-3)}} \\
ReBoot paper, Table 4 (plaintext) & 60k samples & 10 (paper) & 94.61 $\pm$ 0.27 & --- & --- & --- & not reproducible from the released defaults \\
ReBoot released code, paper setting, constant lr, epoch 50 & 60k samples, 50 epochs & 10 & 95.15 $\pm$ 0.19 (n = 10) & --- & 94.61 & +0.54 & reference \\
ReBoot released code, paper setting, constant lr, best epoch (test-selected) & 60k samples, 50 epochs & 10 & 95.35 $\pm$ 0.10 (n = 10) & --- & 94.61 & +0.74 & reference \\
ReBoot released code, paper setting, ReBoot's scheduler, epoch 50 & 60k samples, 50 epochs & 10 & 95.25 $\pm$ 0.08 (n = 10) & --- & 94.61 & +0.64 & reference \\
ReBoot released code, paper setting, ReBoot's scheduler, best epoch (test-selected) & 60k samples, 50 epochs & 10 & 95.46 $\pm$ 0.09 (n = 10) & --- & 94.61 & +0.85 & reference \\
\bottomrule
\end{tabular}}
\par\smallskip\begin{minipage}{\linewidth}\scriptsize
FULL-1 final parameter error: 0.00682657 (maximum absolute difference across weights and biases against its same-batch plaintext twin at training end). This is a trajectory diagnostic, not a per-bootstrap precision or accuracy gate.\par
R-M3-15: the $\pm$0.5 pt gate compares hekit encrypted B = 48 (FULL-2) with ReBoot's released code at its full setting, seed 0. Subset rows are 1 epoch and are not a gate. The paper value and R-REB-3 are reported with this note only.\par
\end{minipage}
\end{table*}

% generated by tools/paper_tables.py from results/; do not edit
\begin{table*}[t]
\centering
\scriptsize
\caption{Plaintext dataset reproduction against ReBoot Table 4 (Task 51)}
\label{tab:plain_datasets}
\resizebox{\linewidth}{!}{%
\begin{tabular}{lrrrrrrrr}
\toprule
Dataset & Config & Epochs & Seeds & Fixed final epoch (\%) & Best epoch, test-selected (\%) & ReBoot Table 4 (\%) & $\Delta$ fixed (pt) & Within 1 pt \\
\midrule
MNIST 14x14 & paper setting & 50 & 10 & 95.00 $\pm$ 0.18 (n = 10) & 95.33 $\pm$ 0.15 (n = 10) & 94.61 $\pm$ 0.27 & +0.39 & yes \\
Fashion-MNIST 14x14 & paper setting & 50 & 10 & 85.22 $\pm$ 0.17 (n = 10) & 85.53 $\pm$ 0.12 (n = 10) & 85.12 $\pm$ 0.20 & +0.10 & yes \\
Breast Cancer & fixed split, min-max (primary) & 100 & 10 & 96.40 $\pm$ 1.27 (n = 10) & 99.30 $\pm$ 0.37 (n = 10) & 99.03 $\pm$ 1.78 & -2.63 & no \\
Breast Cancer & per-run split (sensitivity) & 100 & 10 & 95.26 $\pm$ 1.61 (n = 10) & 97.89 $\pm$ 0.94 (n = 10) & 99.03 $\pm$ 1.78 & -3.77 & no \\
Breast Cancer & z-score (diverges) & 100 & 10 & diverged (10/10 seeds) & --- & 99.03 $\pm$ 1.78 & --- & no \\
\bottomrule
\end{tabular}}
\par\smallskip\begin{minipage}{\linewidth}\scriptsize
Fixed final epoch is the primary figure. The best epoch is selected on the test set and is optimistic; it is shown only because ReBoot does not state its epoch count. Seeds: the plain binary's own seeds (not ReBoot's initial weights).\par
Fashion-MNIST and Breast Cancer are plaintext twins only (no encrypted runs).\par
\end{minipage}
\end{table*}

% generated by tools/paper_tables.py from results/; do not edit
\begin{table*}[t]
\centering
\scriptsize
\caption{Security estimates (lattice-estimator; log2 rop, minimum over attacks)}
\label{tab:security}
\resizebox{\linewidth}{!}{%
\begin{tabular}{lrrrrrr}
\toprule
Parameters & N & log PQ & Secret & BDGL16 & MATZOV & >= 128 (both) \\
\midrule
hekit (a): n16\_boot main key & 2\textasciicircum{}16 & 1732 (1731.998) & dense ternary (P($\pm$1) = 1/4 each, h $\approx$ N/2) & 131.5 & 129.3 & yes \\
hekit (b): SSE key s -> s' (h' = 32) & 2\textasciicircum{}16 & 294 (q\_0$\cdot$P) & sparse ternary, h' = 32 & 133.8 & 134.2 & yes \\
ReBoot main key (R-REB-4: OpenFHE's own chain) & 2\textasciicircum{}16 & 1730 (1729.585) & sparse ternary, h = 192 (97 / 95) & 121.8 & 121.5 & no \\
\bottomrule
\end{tabular}}
\par\smallskip\begin{minipage}{\linewidth}\scriptsize
ReBoot row: the MITM tier (operation counts; about 2\textasciicircum{}104-2\textasciicircum{}110 memory). Without MITM attacks ReBoot is at 127.8 / 126.2 (BDGL16 / MATZOV), also below 128.\par
Source: the summary table of docs/security.md (the estimator outputs are recorded in that doc; no JSON exists). Same lattice-estimator models for every row; m = N; Core-SVP is in the doc for information only. Equal N does not mean equal security: no equal-security claim.\par
\end{minipage}
\end{table*}

Table~\ref{tab:accuracy_vs_reboot} separates the released full-setting reference, subset results, and published paper accuracy. The released defaults and paper-setting reconstruction are different protocols. We do not choose whichever reference makes the result look best. The registered tolerance is tested by FULL-2 against the captured released full-setting seed-zero reference; that encrypted row has not completed. FULL-1 reports its larger-batch accuracy and same-batch twin separately, without substituting either for this gate. Matching a plaintext twin is necessary to diagnose arithmetic fidelity, but is insufficient to establish baseline-equivalent generalization.

Plaintext dataset experiments (Table~\ref{tab:plain_datasets}) provide broader optimization checks. MNIST and Fashion-MNIST reproduce the stated comparison tolerance at the fixed final epoch. Breast Cancer does not: the fixed-epoch min-max result misses its target, and selecting the best test epoch is optimistic leakage from model selection. The z-score setting diverges for the polynomial activation. Fashion-MNIST and Breast Cancer results are plaintext twins, not encrypted training measurements. The local-block plaintext twins reproduce the released eMLP-2/eMLP-3 values for the recorded seed; their encrypted subset rows are pending. The block experiments do not extend the current headline claim to deeper networks.

Table~\ref{tab:security} applies common estimator models to the actual exposed key moduli and secret distributions. The evaluated main and SSE instances of our preset exceed the chosen operation-count target under those models; the released ReBoot preset does not. This is an estimate over the documented attack search, not a security proof. The smallest sparse-secret values use attacks with extreme memory requirements, and even non-MITM minima have substantial or unreported memory demands. These are not practical attack demonstrations. Equal $N$ and a requested library security label do not establish equal security. Results must remain conditional until either a security-different comparison is explicitly accepted or an independently estimated and remeasured alternative baseline is available.

\subsection{RQ3: device operations and ablation}
% generated by tools/paper_tables.py from results/; do not edit
\begin{table*}[t]
\centering
\scriptsize
\caption{GPU vs CPU vs HEonGPU, n16\_train (R-47-2, criterion medians, ms)}
\label{tab:gpu_vs_cpu_ops}
\resizebox{\linewidth}{!}{%
\begin{tabular}{lrrrrrrr}
\toprule
Operation & Level & hekit GPU & CPU 1 thr & CPU 12 thr & GPU speed-up (1 / 12 thr) & HEonGPU (RTX 4090) & GPU / HEonGPU \\
\midrule
add & 31 & 0.111 & 15.674 & 14.611 & 141x / 132x & none published & --- \\
relinearize & 31 & 4.671 & 639.639 & 240.070 & 137x / 51x & 2.582 (l = 36, I = 4 primes per digit) & 1.81x (nominal) \\
mul + relin + rescale & 31 & 5.823 & 725.209 & 337.731 & 125x / 58x & none published & --- \\
rotate(1) & 31 & 3.990 & 622.046 & 223.880 & 156x / 56x & none published & --- \\
hoisted rotate x16 & 31 & skipped: needs 5580 MiB (2377 MiB free) & 3765.068 & 689.471 & --- & none published & --- \\
hoisted rotate x16 (capped keys) & 9 & 4.984 & 1005.331 & 172.870 & 202x / 35x & none published & --- \\
matvec\_enc d = 64 & 31 & skipped: needs 7276 MiB (1929 MiB free) & 12938.440 & 2325.173 & --- & none published & --- \\
matvec\_enc d = 64 (capped keys) & 9 & 16.120 & 2523.223 & 547.762 & 157x / 34x & none published & --- \\
matvec\_enc d = 64 (capped keys) & 5 & 9.509 & 1391.064 & 314.649 & 146x / 33x & none published & --- \\
matvec\_enc x4, 1 stream (sequential) & 5 & 34.994 & --- & --- & --- & none published & --- \\
matvec\_enc x4, 4 streams (untuned R-47-2) & 5 & 41.673 & --- & --- & --- & none published & --- \\
matvec\_enc x4, re-measured after the fix round & 5 & pending (R-47-3) & --- & --- & --- & --- & --- \\
\bottomrule
\end{tabular}}
\par\smallskip\begin{minipage}{\linewidth}\scriptsize
HEonGPU relinearize: Ozcan and Savas, GPU Implementations of Three Different Key-Switching Methods for Homomorphic Encryption Schemes, IACR ePrint 2025/124, Table 3: CKKS relinearization on a GeForce RTX 4090, log N = 16, 43-bit precision, Method II (hybrid), (I, l\textasciitilde{}, l) = (4, 40, 36), level 35. The previously cited Nov 2024 manuscript version date is unverified; ePrint landing metadata alone does not establish that manuscript date. Caveat: not the same parameters as n16\_train (36 vs 32 Q primes, 4 special primes; I = 4 primes per digit corresponds to alpha = 4, not four digits; n16\_train uses dnum 8); the ratio is nominal. Scaling key switching linearly in l gives about 2.30 ms at l = 32 and a ratio of about 2.0x (docs/m4-report.md section 1.1). ePrint 2024/1543 could not be read (HTTP 403), so other HEonGPU operation values are unconfirmed.\par
Capped keys: the input is dropped to level 9 or 5 and the keys are cut to that level (faithful for that level, not the top-level cost); the CPU uses full keys at the same level. The shared RTX 4090 had 2-3 GiB free, so full-key GPU rows that need more were skipped.\par
\end{minipage}
\end{table*}

% generated by tools/paper_tables.py from results/; do not edit
\begin{table*}[t]
\centering
\scriptsize
\caption{RQ3 ablation: one interactive eMLP-1 iteration under each EvalConfig}
\label{tab:rq3_ablation}
\resizebox{\linewidth}{!}{%
\begin{tabular}{lrrrrrrrrr}
\toprule
Source & Backend / level & hoist bsgs lazy & Rotations & Relins & KS/iter & Closed form & Step median (s) & Step / default & Max |$\Delta$W| vs default \\
\midrule
CPU, n16\_train, level 31 &  &  &  &  &  &  & pending (R-48-4) &  &  \\
GPU, n16\_train, level 7 &  &  &  &  &  &  & pending (R-48-2) &  &  \\
CPU, n16\_train, level 7 &  &  &  &  &  &  & pending (R-48-3) &  &  \\
CPU smoke (insecure N = 2\textasciicircum{}12: counts only) & cpu, level max & 1 1 1 & 171 & 62 & 233 & match & n/a (smoke) & n/a (smoke) & 0 \\
CPU smoke (insecure N = 2\textasciicircum{}12: counts only) & cpu, level max & 1 1 0 & 171 & 114 & 285 & match & n/a (smoke) & n/a (smoke) & 7.83e-09 \\
CPU smoke (insecure N = 2\textasciicircum{}12: counts only) & cpu, level max & 1 0 1 & 210 & 101 & 311 & match & n/a (smoke) & n/a (smoke) & 6.90e-09 \\
CPU smoke (insecure N = 2\textasciicircum{}12: counts only) & cpu, level max & 1 0 0 & 210 & 114 & 324 & match & n/a (smoke) & n/a (smoke) & 7.39e-09 \\
CPU smoke (insecure N = 2\textasciicircum{}12: counts only) & cpu, level max & 0 1 1 & 171 & 62 & 233 & match & n/a (smoke) & n/a (smoke) & 0 \\
CPU smoke (insecure N = 2\textasciicircum{}12: counts only) & cpu, level max & 0 1 0 & 171 & 114 & 285 & match & n/a (smoke) & n/a (smoke) & 7.83e-09 \\
CPU smoke (insecure N = 2\textasciicircum{}12: counts only) & cpu, level max & 0 0 1 & 210 & 101 & 311 & match & n/a (smoke) & n/a (smoke) & 6.90e-09 \\
CPU smoke (insecure N = 2\textasciicircum{}12: counts only) & cpu, level max & 0 0 0 & 210 & 114 & 324 & match & n/a (smoke) & n/a (smoke) & 7.39e-09 \\
\bottomrule
\end{tabular}}
\par\smallskip\begin{minipage}{\linewidth}\scriptsize
1. Definition of bsgs = false. The giant rotation is not shared; the weights keep the same pre-rotated (BSGS) layout and the same Galois keys. Every product with a nonzero giant shift is finished (relinearized) and rotated on its own; the products of the zero-shift block are summed as with BSGS. Per matvec of dimension d (offset 0): d - 1 rotations, of which only the baby - 1 rotations of x can be hoisted; with lazy\_relin 1 + d - baby relins, without it d (EvalConfig::matvec\_rotations, EvalConfig::matvec\_enc\_relins). At d = 16 (baby = giant = 4), as (hoistable rotations / individual rotations / relins): | form | hoistable rot. | individual rot. | relins (lazy) | relins (not lazy) | |---|---|---|---|---| | default BSGS | 3 | 3 | 4 | 16 | | bsgs = false (this library) | 3 | 12 | 13 | 16 | | textbook diagonal method (not implemented) | 15 | 0 | 1 | 16 | The textbook (Halevi-Shoup) diagonal method rotates x by all d - 1 steps (all hoistable) and sums the products lazily; it needs un-pre-rotated diagonals, i.e. another weight layout and gradient form, so it is not implemented. The all-off configuration has the same counts in both forms (15 rotations, 16 relins at d = 16).\par
2. Bias. Against the strongest naive baseline (the textbook method with hoisting and lazy relinearization), the cost of the bsgs = false rows overstates what BSGS gains: the measured BSGS gain is an upper bound. Conversely, the marginal gains of hoist and lazy\_relin measured inside bsgs = false configurations are lower bounds (only baby - 1 rotations can be hoisted and only the zero-shift block summed lazily: e.g. lazy\_relin saves 3 relins at d = 16 instead of 15).\par
3. hoist changes no OpCounter value. The counter does not distinguish hoisted from individual rotations, so the effect of hoist shows only in wall-clock time, never in key-switch counts.\par
4. bsgs does not touch the gradient's outer product (Dense::weight\_grad\_scaled, Dense::weight\_grad\_reduced): its (baby - 1) + (giant - 1) input rotations are what the pre-rotated weight layout requires, not a BSGS choice; only hoist changes how they run. The interactive gradient multiplies each diagonal with one Evaluator::mul (nothing to sum lazily), so lazy\_relin does not change it either.\par
The smoke rows (insecure N = 2\textasciicircum{}12, high Mac load) check counts and closed forms only; their times are not comparative.\par
\end{minipage}
\end{table*}

Table~\ref{tab:gpu_vs_cpu_ops} contains synchronized operation latencies, including capped-level rows and memory skips. The HEonGPU relinearization reference~\cite{gpuKeys2025} is nominal: its modulus count differs. A linear-in-limb adjustment is a cost-model extrapolation, not an equal-parameter measurement. The published table value is inherited from the archived operator report; the bibliographic metadata was checked for this draft, but a fresh PDF fetch was blocked. The later pool/submission fix has not been retimed.

Table~\ref{tab:rq3_ablation} deliberately leaves the server rows pending. The insecure small-parameter CPU smoke validates counts and checks only; its high-load timings are inadmissible for performance comparisons. The flag named \texttt{bsgs=false} keeps pre-rotated weights rather than implementing the strongest textbook diagonal baseline. Thus its BSGS gain would be an upper bound against that alternative, while the marginal hoisting and lazy-relinearization gains within those rows would be lower bounds. Hoisting changes no operation counter; the gradient outer product retains layout-required rotations independent of the BSGS flag. All four caveats travel with the table.

CPU and GPU training comparisons require the same optimizer, level, layout, and batch. Registered SGD fallback rows are compared only with the corresponding CPU SGD row. A resource skip is unmeasured; a device OOM is a failure. Test-process success cannot stand in for an actually executed production-size parity gate. GPU bootstrap precision, latency, allocation behavior, and peak memory are still unmeasured. HEonGPU's separate bootstrap reference~\cite{heongpuBench} has a different remaining-level budget and main-secret distribution; the fixed plan target is not a normalized speed-ratio result. RQ3's epoch-time question is therefore unanswered.

\section{Related Work, Limitations, and Conclusion}
\label{sec:limits}
\subsection{Positioning}
The earlier diagonal-training design~\cite{mihara2020} motivates rotation-only transposition. We generalize the layout and charge server-side batch maintenance, compact refresh, and optimizer state. ReBoot~\cite{reboot2026} already establishes non-interactive training and local learning; its released implementation is our application baseline, not a source of imported implementation code. HETAL~\cite{hetal2023} studies encrypted transfer learning and efficient matrix products in a different training scope. Our evaluation does not reproduce its workloads or establish a ranking against it.

The cryptographic building blocks are efficient linear transforms and hoisting~\cite{halevi2018}, SSE~\cite{sse2022}, and CKKS~\cite{ckks2017}. SEAL~\cite{seal} provides interoperability checks; OpenFHE~\cite{openfhe2022} supplies the separately executed baseline. HEonGPU~\cite{gpuKeys2025,heongpuBench} is a device-operation reference, not an end-to-end training competitor measured under our protocol. An MIT implementation improves reuse, but licensing and pipeline completeness are not evidence of a cryptographic or learning-performance advance. The approved dependency policy permits MIT, Apache, BSD, and ISC licenses, with one narrow Unicode-3.0 character-table exception for \texttt{unicode-ident} version \texttt{1.0.27}. That component requires its complete Unicode notice; this is not a general Unicode-license exception.

\subsection{Limits and missed goals}
\paragraph{Scope and comparison.}
FULL-1 and FULL-2 supply completed full-setting encrypted runs at two batches with the same initial-weight seed, including their schedules and accumulated error. FULL-2 passes the registered baseline-batch accuracy gate, but misses the common-baseline-bootstrap count target. FULL-1's successful count comparisons are a separate configuration. Full-setting ReBoot CPU wall-clock is unmeasured, and the common-security condition is unmet. Neither these two configurations nor the earlier subsets establish multi-seed encrypted behavior; multi-seed plaintext results are a different evidence class. The observed parameter error does not establish a general error bound. Batch choices trade refresh amortization against learning dynamics. The effective baseline security differs, initialization and bias arithmetic require care, and legacy rows cannot support accuracy claims. ReBoot's in-library layout timing was measured on a different machine and is not a same-CPU baseline; its public-label/client-update path has a weaker confidentiality contract. The official encrypted block-model baselines have not been run: their reported memory demand exceeds the available host budget.

\paragraph{Cost and correctness.}
A key switch at a low level need not cost the same as one at the top level; counts omit memory traffic, cache effects, and key residency. The initial compact-layout assumptions required actual CKKS tests and the first model was optimistic. Replicated copies were once unstable during interactive training; canonical-layout refresh was required to avoid divergence. The earlier paper's published timing total and phase sum disagree, so the continuity demonstration does not silently reconcile them or infer a server-only speedup. The original CPU NTT target and several SEAL-relative key-switch performance targets were missed. The early CUDA NTT spike was not timed against its registered hardware reference; later device-operator measurements do not retroactively supply that measurement. Some SEAL comparisons also have different special-prime/key-switch decompositions. We report these as missed performance targets, not failed numerical gates.

\paragraph{Numerical limits.}
SSE and modulus-band changes were needed because the planned dense-secret overflow range and original prime allocation did not meet the bootstrap precision budget. Unit-scale nonconstant masks were invalid, and exact-scale folding is required to fit the training circuit. The accepted security claim uses specified sieve cost models, not a universal security level; Core-SVP estimates are lower. Dense-secret hybrid search is bounded rather than an exhaustive global optimization. An all-sparse main-key alternative did not meet the selected security criterion and is not used. KAT coverage is limited to the supported SEAL preset, with asymmetric comparison directions; it does not independently validate our bootstrap or larger hybrid preset. Standalone OpenFHE/Lattigo bootstrap precision parity remains absent. Random full-slot CPU precision checks do not prove a negligible failure probability for arbitrary training distributions or activations outside range.

\paragraph{Learning limits.}
The quartic activation's claimed accuracy benefit is unmeasured. The evaluated encrypted trainer omits weight decay and optional softmax/cross-entropy. The separate plaintext paper-setting reconstruction includes weight decay; it is not the faithful encrypted protocol. Bias momentum is not matched to NAG, and faithful bias freezing preserves extra encrypted arithmetic. Fixed-final-epoch Breast Cancer accuracy misses its target; test-selected best epochs are not a remedy. No encrypted Fashion-MNIST or Breast Cancer accuracy is reported. Higher-depth local-block training has plaintext fidelity checks but no completed encrypted comparison. The continuity Iris experiment is interactive and differs in padding, randomization, and normalization; it cannot establish the primary non-interactive claim. Threshold or multi-party FHE, scheme switching, transformer training, and Metal execution are outside the evaluated scope.

\paragraph{Device and implementation limits.}
Full-key GPU training, the matched training parity gate, the RQ3 timing ablation, and GPU bootstrap measurements are pending. Repaired multi-stream latency is measured only for independent products at a capped level; it cannot supply those missing training results. Memory estimates and pool snapshots do not establish a peak. All-off flags do not implement the strongest naive diagonal algorithm. The current training benchmark still uses host bootstrap fallback. Device error handling and numerical replay are useful engineering evidence, but compile-only paths and skipped tests are not physical hardware results. CPU/GPU side-channel resistance, constant-time sampling, secret zeroization, malformed-input hardening, and independent cryptographic audit are outside the current implementation's assurance. Some internal invariants are debug assertions; decode/parameter mismatches can panic. The artifact is research software, not a deployment security claim.

\subsection{Conclusion}
Batched diagonal packing preserves a rotation-only transpose while supporting compact non-interactive optimizer updates. The completed full-setting runs show a larger-batch count advantage and baseline-batch accuracy within tolerance, together with a missed normalized count target at the baseline batch. The security mismatch and unmeasured full-setting baseline time prevent an unconditional superiority claim. The device microbenchmark also shows that stream and pool choices matter, without resolving training performance. The resulting evidence supports a conditional re-evaluation of the layout and identifies the matched-security, baseline-timing, and hardware experiments still needed.

\bibliographystyle{ACM-Reference-Format}
\bibliography{refs}
\appendix
\section{Parameters, Security, and Reproduction}
\label{app:params}
\subsection{Implemented chain and bootstrap contract}
The production bootstrap preset uses $N=2^{16}$, $N/2$ real slots, scale $2^{42}$, and Q-prime bit sizes
\begin{equation*}
 [54,\underbrace{42,\ldots,42}_{16},\underbrace{34,\ldots,34}_{4},
 \underbrace{58,\ldots,58}_{7},\underbrace{56,\ldots,56}_{4}].
\end{equation*}
Four special primes have bit size $60$, with eight hybrid digits. The bands allocate four levels to each linear transform and seven to EvalMod, leaving level $16$ before training unpack. EvalMod uses Chebyshev degree $127$, interval parameter $K=13$, and no double-angle step. These are implemented design constants; the superseded original-plan bands and dense-secret range are not used.

The bootstrap requires level zero and exactly the preset scale before work. Dense-to-sparse SSE is restricted to $q_0P$ under a temporary key of weight $32$; the return key and remaining arithmetic use the dense main key over divisors of $QP$. CPU full-slot tests exceed the registered precision threshold for random real inputs in $[-1,1]$. This does not imply a worst-case failure bound. The GPU implementation has local replay and compilation coverage, but its physical precision and memory feasibility remain unmeasured.

\subsection{What the estimator models}
Our default sampler gives $-1,0,1$ probabilities $1/4,1/2,1/4$. We model the dense secret with a fixed-count sparse-ternary proxy rather than a uniform-ternary distribution; SSE is a separate small-modulus instance. The reported attack input rounds the exposed $\log_2(QP)$ upward and uses $m=N$ samples. Dense and sparse main keys must not be interchanged at the same modulus without re-estimation.

The estimator is pinned to commit \texttt{ae45d740ceb1}; methodology follows Albrecht et al.~\cite{albrecht2015} and the maintained estimator documentation~\cite{latticeEstimator}. The recorded search includes primal/dual and selected hybrid/MITM variants. Dense primal-hybrid optimization did not finish globally; fixed guessing-dimension scans supplement it. Arora--GB and BKW were not included. For the released ReBoot context, an observed sign split from one generated key is an input instance, not a distributional guarantee for all keys. The minima are conditional on the selected search, cost model, and memory assumptions. Table~\ref{tab:security} is not a proof, nor evidence of a practical compromise.

\subsection{Transpose identity and reduction boundary}
For Eq.~\eqref{eq:transpose}, matching row and column indices gives $t\equiv i\pmod {n_o}$ and $t+c\equiv0\pmod {n_i}$. The wide/tall cases in Section~\ref{sec:packing} solve both conditions. Substituting the pre-rotations on the source and destination yields the displayed offset. Because every weight diagonal is periodic with period $L$ in feature coordinates and all offsets are multiples of $B$, samples are not mixed. The identity therefore applies to rectangular padded layers and repeated blocks as well as square matrices. A zero offset is a copy; other offsets require the selected dedicated or composed rotation keys.

Lane reduction is different: offsets not divisible by $B$ mix lanes intentionally, and isolating one lane requires a mask. Ignoring that mask or its replication work changes both the depth and count model. A partial-batch validity mask must also exclude unused lanes from the loss derivative; zero input alone is insufficient when the model has a bias or a linear output loss.

\subsection{KAT and artifact boundaries}
The SEAL compatibility checks compare encoding, serialized ciphertext interoperability, decryption, and supported evaluation paths under the shared preset. Some directions assert exact residues while others only assert a decoded-error bound. Rotation arithmetic intentionally differs in its centered lifting. Neither this KAT nor CPU/GPU bit equality establishes security; the latter only checks implementation agreement for executed cases. Small insecure presets are smoke tests and never enter the security claim.

Tables are generated by \texttt{tools/paper\_tables.py}; field-level provenance is available through its \texttt{--sources} option. The figure generator consumes those validated table cells and writes a source manifest containing paths, fields, and SHA-256 digests. It does not rewrite result JSON. Missing registered runs produce explicit pending rows; existing invalid inputs must fail validation. Timed runs, setup, client diagnostics, and model-only extrapolations stay distinguishable. Running \texttt{scripts/run\_all.sh --tables-only} regenerates the retained tables. This does not execute the complete registered experiment matrix, whose end-to-end runner remains unfinished; a quick pipeline check is not a source of research results.

The anonymous review mode omits author metadata and identifying artifact links. ReBoot and OpenFHE are executed as external baselines; no implementation source is copied into this artifact. The estimator and TeX compiler are isolated development tools, not library dependencies. The library's dependency policy, including the specific character-table exception disclosed in Section~\ref{sec:limits}, is separate from the licenses of those external research tools. Instructions for rebuilding anonymous review, named review, and named preprint variants are provided with the source. Authorship is confirmed for this preprint. Before arXiv submission, the final source, metadata, and server-generated PDF require checking against the retained evidence; unavailable measurements remain unavailable. A later conference submission requires a separate anonymity and venue-rules review; the WAHC~2026 limit is only the provisional formatting constraint~\cite{wahc2026}.

\ifdefined\ArxivDraft
Code and evidence: \url{https://github.com/kenmaro3/ckks-train}. Release \texttt{v1.0.0-preview.1} is an earlier snapshot; this preprint has its own source revision.
\fi

\end{document}
